\chapter{Números reais} \label{rn:chapter}

%%%%%%%%%%%%%%%%%%%%%%%%%%%%%%%%%%%%%%%%%%%%%%%%%%%%%%%%%%%%%%%%%%%%%%%%%%%%%%

\section{Propriedades básicas} \label{sec:basicpropsrn}

\sectionnotes{1,5 aulas}

Em análise, o objeto principal com que trabalhamos é o conjunto dos
\myindex{números reais}. Como esse conjunto é tão fundamental, muitas vezes
dedica-se bastante tempo à sua construção formal. No entanto, adotaremos uma
abordagem mais simples e suporemos que existe um conjunto com as propriedades
adequadas.
As três propriedades fundamentais dos números reais são: ser um conjunto
ordenado, ser completo em relação a essa ordem e ser um corpo compatível com
essa ordem.
Começamos pela ordem.

\begin{defn}
Um \emph{\myindex{conjunto ordenado}} é um conjunto $S$ munido de
uma relação $<$ tal que
%\begin{enumerate}[(i),itemsep=0.5\itemsep,parsep=0.5\parsep,topsep=0.5\topsep,partopsep=0.5\partopsep]
%\begin{enumerate}[(i),nolistsep]
\begin{enumerate}[(i)]
\item \emph{(\myindex{tricotomia})} Para todos $x, y \in S$, exatamente uma
das afirmações $x < y$, $x=y$ ou $y < x$ é verdadeira.
\item \emph{(transitividade)} Se $x,y,z \in S$ são tais que $x < y$ e $y
< z$, então $x < z$.
\end{enumerate}
Escrevemos $x \leq y$ se $x < y$ ou $x=y$. Definimos
$>$ e $\geq$ da maneira usual.
\glsadd{not:ordereq}
\glsadd{not:orderlt}
\glsadd{not:orderleq}
\glsadd{not:ordergt}
\glsadd{not:ordergeq}
\end{defn}


O conjunto dos números racionais $\Q$ é ordenado: dizemos
$x < y$ se, e somente se, $y-x$ é um número racional positivo, isto é,
$y-x = \nicefrac{p}{q}$, onde $p,q \in \N$. Da mesma forma,
$\N$ e $\Z$ também são conjuntos ordenados.

Há conjuntos ordenados que não são conjuntos de números. Por exemplo, o
conjunto dos países pode ser ordenado por área territorial, de modo que a Índia $>$ Liechtenstein.
Um conjunto ordenado típico que você conhece desde o ensino fundamental é o
dicionário. Ele é o conjunto ordenado das palavras, cuja ordem é chamada de
ordem lexicográfica. Conjuntos ordenados desse tipo aparecem com frequência,
por exemplo, na ciência da computação. Neste livro, nosso interesse estará
principalmente nos conjuntos ordenados de números.

\begin{defn}
Seja $E \subset S$, onde $S$ é um conjunto ordenado.
%\begin{enumerate}[(i),itemsep=0.5\itemsep,parsep=0.5\parsep,topsep=0.5\topsep,partopsep=0.5\partopsep]
%\begin{enumerate}[(i),nolistsep]
\begin{enumerate}[(i)]
\item Se existe $b \in S$ tal que $x \leq b$ para todo $x \in E$,
dizemos que $E$ é \emph{\myindex{limitado superiormente}} e que $b$
é uma \emph{\myindex{cota superior}} de $E$.
\item Se existe $b \in S$ tal que $x \geq b$ para todo $x \in E$,
dizemos que $E$ é \emph{\myindex{limitado inferiormente}} e que $b$
é uma \emph{\myindex{cota inferior}} de $E$.
\item Se existe uma cota superior $b_0$ de $E$ tal que
$b_0 \leq b$ para toda cota superior $b$ de $E$, então
$b_0$ é chamada de \emph{\myindex{menor cota superior}} ou
\emph{\myindex{supremo}}
de $E$. Veja a \figureref{fig:lub}. Escrevemos\glsadd{not:sup}
\begin{equation*}
\sup\, E \coloneqq b_0 .
\end{equation*}
\item De maneira semelhante, se existe uma cota inferior $b_0$ de $E$ tal que
$b_0 \geq b$ para toda cota inferior $b$ de $E$, então
$b_0$ é chamada de \emph{\myindex{maior cota inferior}} ou
\emph{\myindex{ínfimo}}
de $E$. Escrevemos\glsadd{not:inf}
\begin{equation*}
\inf\, E \coloneqq b_0 .
\end{equation*}
\end{enumerate}
Quando um conjunto $E$ é limitado tanto superior quanto inferiormente,
dizemos simplesmente que $E$ é \emph{limitado}\index{conjunto limitado}.
\end{defn}

A notação $\sup E$ e $\inf E$ é justificada porque o
supremo (ou ínfimo) é único (quando existe): Se $b$ e
$b'$ são supremos de $E$, então $b \leq b'$ e $b' \leq b$, pois $b$ e $b'$ são ambos menores cotas superiores; portanto, $b=b'$.


\begin{myfigureht}
\myincludepdft{lub}{%
Um diagrama do conjunto E na reta real, representado por uma linha horizontal.
Uma seta à esquerda indica que os números diminuem quando nos deslocamos para
a esquerda, e outra seta à direita indica que os números aumentam quando nos
deslocamos para a direita. O conjunto E aparece em trechos destacados da reta
real. Vários traços verticais, todos à direita do conjunto E, estão marcados
como cotas superiores de E; o menor deles, situado exatamente na extremidade
de E, está marcado como a menor cota superior de E.}
\caption{Um conjunto $E$ limitado superiormente e a menor cota superior de $E$.\label{fig:lub}}
\end{myfigureht}

Um exemplo simples:
Seja $S \coloneqq \{ a, b, c, d, e \}$ ordenado por $a < b < c < d < e$, e
seja $E \coloneqq \{ a, c \}$. Então $c$, $d$ e $e$ são cotas superiores de $E$, e
$c$ é a menor cota superior, ou supremo, de $E$.

O supremo ou o ínfimo de $E$ (mesmo quando existe) não precisa pertencer a
$E$. O conjunto $E \coloneqq \{ x \in \Q : x < 1 \}$ tem 1 como menor cota
superior, mas 1 não pertence ao próprio conjunto $E$.
O conjunto
$G \coloneqq \{ x \in \Q : x \leq 1 \}$ também tem 1 como cota superior,
e, neste caso, $1 \in G$.
O conjunto $P \coloneqq \{ x \in \Q : x \geq 0 \}$ não tem cota superior (por quê?) e, portanto,
não pode ter menor cota superior. O conjunto $P$ tem uma maior cota
inferior:~0.

\begin{defn} \label{defn:lub}
Um conjunto ordenado $S$ tem a \emph{\myindex{propriedade da menor cota superior}} se
todo subconjunto não vazio
$E \subset S$ limitado superiormente tem uma menor cota superior, isto é,
$\sup\, E$ existe em $S$.
\end{defn}

A \emph{propriedade da menor cota superior}
também é chamada, às vezes, de \emph{\myindex{propriedade da completude}} ou
\emph{\myindex{propriedade da completude de Dedekind}}%
\footnote{%
Em homenagem ao matemático alemão
\href{https://en.wikipedia.org/wiki/Richard_Dedekind}{Julius Wilhelm Richard Dedekind}
(1831--1916).}.
Os números reais têm essa propriedade.

\begin{example}
O conjunto $\Q$ dos números racionais não tem a propriedade da menor cota superior. O subconjunto
$\{ x \in \Q : x^2 < 2 \}$ não tem supremo em $\Q$. Veremos mais adiante (\exampleref{example:sqrt2}) que o supremo é
$\sqrt{2}$, que não é racional%
\footnote{O mesmo vale para todas as outras raízes de 2; além disso, o fato de
$\sqrt[k]{2}$ nunca ser racional para $k > 1$ está relacionado à
impossibilidade de afinar perfeitamente um piano em todas as tonalidades. Veja, por exemplo:
\url{https://youtu.be/1Hqm0dYKUx4}.}.
Suponha que $x \in \Q$ seja tal que $x^2 = 2$.
Escreva $x=\nicefrac{m}{n}$ em forma irredutível. Então ${(\nicefrac{m}{n})}^2 = 2$
ou
$m^2 = 2n^2$. Portanto, $m^2$ é divisível por 2, e assim $m$ é divisível por
2. Escreva $m = 2k$; então ${(2k)}^2 = 2n^2$. Divida por 2
e observe que $2k^2 = n^2$; logo, $n$ é divisível por 2. Mas isso é uma
contradição, pois $\nicefrac{m}{n}$ está em forma irredutível.
\end{example}

O fato de $\Q$ não ter a propriedade da menor cota superior é uma das
principais razões pelas quais trabalhamos com $\R$ em análise. O conjunto $\Q$
é perfeitamente adequado para os algebristas. Mas nós, analistas, precisamos
da propriedade da menor cota superior para fazer qualquer coisa.
Também exigimos que os números reais tenham muitas propriedades algébricas.
Em particular, exigimos que formem um corpo.


%\enlargethispage{\baselineskip}
\begin{defn}
Um conjunto $F$ é chamado de \emph{\myindex{corpo}} se nele estiverem
definidas duas operações, a adição $x+y$ e a multiplicação $xy$, e se forem
satisfeitos os seguintes axiomas\index{axioma}:
%mbxSTARTIGNORE
\begin{enumerate}[({A}1)]
%mbxENDIGNORE
%mbxCLOSEPARAGRAPH
%mbx <dl>
\item
%mbx <title>(A1)</title>
Se $x \in F$ e $y \in F$, então $x+y \in F$.
\item
%mbx <title>(A2)</title>
\emph{(comutatividade da adição)}
$x+y = y+x$ para todos $x,y \in F$.
\item
%mbx <title>(A3)</title>
\emph{(associatividade da adição)}
$(x+y)+z = x+(y+z)$ para todos $x,y,z \in F$.
\item
%mbx <title>(A4)</title>
Existe um elemento $0 \in F$ tal que
$0+x = x$ para todo $x \in F$.
\item
%mbx <title>(A5)</title>
Para cada elemento $x\in F$, existe um elemento $-x \in F$
tal que $x + (-x) = 0$.
%mbxCLOSEITEM
%mbx </dl>
%mbxSTARTIGNORE
\end{enumerate}
\begin{enumerate}[({M}1)]
%mbxENDIGNORE
%mbx <dl>
\item
%mbx <title>(M1)</title>
Se $x \in F$ e $y \in F$, então $xy \in F$.
\item
%mbx <title>(M2)</title>
\emph{(comutatividade da multiplicação)}
$xy = yx$ para todos $x,y \in F$.
\item
%mbx <title>(M3)</title>
\emph{(associatividade da multiplicação)}
$(xy)z = x(yz)$ para todos $x,y,z \in F$.
\item
%mbx <title>(M4)</title>
Existe um elemento $1 \in F$ (com $1 \neq 0$) tal que
$1x = x$ para todo $x \in F$.
\item
%mbx <title>(M5)</title>
Para cada $x\in F$ tal que $x \neq 0$, existe um elemento
$\nicefrac{1}{x} \in F$
tal que $x(\nicefrac{1}{x}) = 1$.
%\end{enumerate}
%\begin{enumerate}
%mbxSTARTIGNORE
\item[(D)]
%mbxENDIGNORE
%mbxlatex \item
%mbx <title>(D)</title>
\emph{(lei distributiva)} $x(y+z) = xy+xz$
para todos $x,y,z \in F$.
%mbxSTARTIGNORE
\end{enumerate}
%mbxENDIGNORE
%mbxCLOSEITEM
%mbx </dl>
\end{defn}

\begin{example}
O conjunto $\Q$ dos números racionais é um corpo. Por outro lado, $\Z$ não é um
corpo, pois não contém os inversos multiplicativos. Por exemplo, não existe
$x \in \Z$ tal que $2x = 1$, portanto o axioma (M5) não é satisfeito. Você
pode verificar que (M5) é a única propriedade que falha\footnote{Um algebrista diria que $\Z$ é um anel ordenado ou, talvez, um anel comutativo ordenado.}.
\end{example}

Vamos supor conhecidos os fatos básicos sobre corpos que se demonstram
facilmente a partir dos axiomas. Por exemplo, é fácil demonstrar que $0x = 0$,
observando que $xx = (0+x)x = 0x+xx$, usando (A4), (D)\@ e (M2). Então,
usando (A5) em $xx$, junto com (A2), (A3) e (A4), obtemos $0 = 0x$.

\begin{defn}
Dizemos que um corpo $F$ é um \emph{\myindex{corpo ordenado}} se
$F$ também é um conjunto ordenado tal que
\begin{enumerate}[(i)]
\item \label{defn:ordfield:i} Para $x,y,z \in F$, $x < y$ implica $x+z <
y+z$.
\item \label{defn:ordfield:ii} Para $x,y \in F$, $x > 0$ e $y > 0$
implicam $xy > 0$.
\end{enumerate}
Se $x > 0$, dizemos que $x$ é \emph{\myindex{positivo}}.
Se $x < 0$, dizemos que $x$ é \emph{\myindex{negativo}}.
Também dizemos que $x$ é \emph{\myindex{não negativo}} se $x \geq 0$
e que $x$ é \emph{\myindex{não positivo}} se $x \leq 0$.
\end{defn}

Os números racionais $\Q$ com a ordem usual formam um corpo ordenado.
Deixamos os detalhes para o leitor interessado.

\begin{prop} \label{prop:bordfield}
Seja $F$ um corpo ordenado e sejam $x,y,z,w \in F$. Então
\begin{enumerate}[(i)]
\item \label{prop:bordfield:i} Se $x > 0$, então $-x < 0$ (e vice-versa).
\item \label{prop:bordfield:ii} Se $x > 0$ e $y < z$, então $xy < xz$.
\item \label{prop:bordfield:iii} Se $x < 0$ e $y < z$, então $xy > xz$.
\item \label{prop:bordfield:iv} Se $x \neq 0$, então $x^2 > 0$.
\item \label{prop:bordfield:v} Se $0 < x < y$, então $0 < \nicefrac{1}{y} < \nicefrac{1}{x}$.
\item \label{prop:bordfield:vi} Se $0 < x < y$, então $x^2 < y^2$.
\item \label{prop:bordfield:vii} Se $x \leq y$ e $z \leq w$, então $x + z \leq y + w$.
\end{enumerate}
\end{prop}

Observe que \ref{prop:bordfield:iv} implica, em particular, que $1 > 0$.

\begin{proof}
Vamos provar \ref{prop:bordfield:i}. A desigualdade $x > 0$ implica, pelo item
\ref{defn:ordfield:i} da definição de corpo ordenado, que
$x + (-x) > 0 + (-x)$. Usando as propriedades algébricas dos corpos,
obtemos $0 > -x$. A afirmação \myquote{e vice-versa} segue por um cálculo
semelhante.

Para provar \ref{prop:bordfield:ii}, observe que $y < z$ implica
$0 < z - y$ pelo
item \ref{defn:ordfield:i} da definição de corpo ordenado.
Aplique o item
\ref{defn:ordfield:ii} da definição de corpo ordenado para obter
$0 < x(z-y)$. Pelas propriedades algébricas, $0 < xz - xy$.
Novamente pelo item \ref{defn:ordfield:i} da definição, $xy < xz$.

A parte \ref{prop:bordfield:iii} fica como exercício.

Para provar a parte \ref{prop:bordfield:iv}, suponha primeiro que $x > 0$.
Pelo item \ref{defn:ordfield:ii} da definição de corpo ordenado,
$x^2 > 0$ (basta tomar $y=x$). Se $x < 0$, usamos
a parte \ref{prop:bordfield:iii} desta proposição, tomando $y=x$ e
$z=0$.

Para provar a parte \ref{prop:bordfield:v}, observe que
$\nicefrac{1}{x}$ não pode ser igual a zero (por quê?).
Suponha que $\nicefrac{1}{x} < 0$;
então $\nicefrac{-1}{x} > 0$ pela parte \ref{prop:bordfield:i}.
Aplique a parte \ref{defn:ordfield:ii} da definição (pois $x > 0$) para obter
$x(\nicefrac{-1}{x}) > 0$ ou $-1 > 0$, o que contradiz $1 > 0$ usando novamente a parte
\ref{prop:bordfield:i}. Portanto, $\nicefrac{1}{x} > 0$.
De maneira semelhante, $\nicefrac{1}{y} > 0$. Logo $(\nicefrac{1}{x})(\nicefrac{1}{y}) > 0$
pela definição de corpo ordenado e, pela parte \ref{prop:bordfield:ii},
\begin{equation*}
(\nicefrac{1}{x})(\nicefrac{1}{y})x < (\nicefrac{1}{x})(\nicefrac{1}{y})y .
\end{equation*}
Pelas propriedades algébricas, $\nicefrac{1}{y} < \nicefrac{1}{x}$.

As partes \ref{prop:bordfield:vi} e \ref{prop:bordfield:vii} ficam como
exercícios.
\end{proof}

O produto de dois números positivos (elementos de um corpo ordenado) é positivo
pela definição de corpo ordenado.
No entanto, não é verdade que, se o produto é positivo, cada um dos fatores
precisa ser positivo. Por exemplo, $(-1)(-1) = 1 > 0$.

\begin{prop}
Sejam $x,y \in F$, onde $F$ é um corpo ordenado. Se
$xy > 0$, então ou ambos $x$ e $y$ são positivos, ou ambos são negativos.
\end{prop}

\begin{proof}
Vamos demonstrar a contrapositiva: \emph{Se $x$ ou $y$ é zero, ou
se $x$ e $y$ têm sinais opostos, então $xy$ não é positivo.}
Se $x$ ou $y$ é zero, então $xy$ é zero e, portanto, não é positivo.
Assim, suponha que $x$ e $y$ sejam não nulos e tenham
sinais opostos.
Sem perda de generalidade,
suponha que $x > 0$ e $y < 0$.
Multiplique $y < 0$ por $x$ para obter $xy < 0x = 0$.
\end{proof}

\begin{example} \label{example:complexfield}
Talvez o leitor já tenha ouvido falar dos \emph{\myindex{números complexos}},
geralmente denotados por $\C$.\glsadd{not:complex}
Isto é, $\C$ é o conjunto dos números da forma
$x + iy$, onde $x$ e $y$ são números reais e $i$ é a
unidade imaginária, um número tal que $i^2 = -1$. Talvez o leitor
se lembre, das aulas de álgebra, de que $\C$ também é um corpo; no entanto, não é um
corpo ordenado. Embora seja possível tornar $\C$ um conjunto ordenado de
alguma maneira, não é possível definir em $\C$ uma ordem que o torne um
corpo ordenado: Em todo corpo ordenado,
$-1 < 0$ e $x^2 > 0$ para todo $x$ não nulo, mas em $\C$ temos
$i^2 = -1$.
\end{example}

Por fim, um corpo ordenado que tem a propriedade da menor cota superior também
tem a propriedade correspondente para as maiores cotas inferiores.

\begin{prop}
Seja $F$ um corpo ordenado com a propriedade da menor cota superior.
Seja $A \subset F$ um conjunto não vazio e limitado inferiormente.
Então $\inf\, A$ existe.
\end{prop}

\begin{proof}
Seja $B \coloneqq \{ -x : x \in A \}$. Seja $b \in F$ uma cota inferior de $A$:
Se $x \in A$, então $x \geq b$
e, portanto, $-x \leq -b$. Assim, $-b$
é uma cota superior de $B$.
Como $F$ tem a propriedade da menor cota superior, existe $c\coloneqq\sup\, B$, e $c \leq -b$.
Como $y \leq c$
para todo $y \in B$, temos $-c \leq x$ para todo $x \in A$.
Logo,
$-c$ é uma cota inferior de $A$. Como $-c \geq b$
e $b$ é uma cota inferior arbitrária de $A$,
a maior cota inferior de $A$ existe e é igual a $-c$.
\end{proof}

\subsection{Exercícios}

\begin{exercise}
Seja $F$ um corpo ordenado e sejam $x,y,z \in F$. Demonstre a parte
\ref{prop:bordfield:iii} de \propref{prop:bordfield}.
Isto é, demonstre que, se $x < 0$ e $y < z$, então $xy > xz$.
\end{exercise}


\begin{exercise} \label{exercise:finitesethasminmax}
Seja $S$ um conjunto ordenado. Seja $A \subset S$ um subconjunto finito e não vazio.
Mostre que $A$ é limitado. Além disso,
mostre que
$\inf\, A$ existe e pertence a $A$ e que
$\sup\, A$ existe e pertence a $A$.
Dica: use a \hyperref[induction:thm]{indução}.
\end{exercise}

\begin{exercise} \label{exercise:squareineq}
Demonstre a parte \ref{prop:bordfield:vi} de \propref{prop:bordfield}.
Isto é, sejam $x, y \in F$, onde $F$ é um corpo ordenado, tais que
$0 < x < y$. Mostre que $x^2 < y^2$.
\end{exercise}

\begin{exercise}
Seja $S$ um conjunto ordenado. Seja $B \subset S$ um conjunto limitado (superior
e inferiormente). Seja $A \subset B$ um subconjunto não vazio.
Suponha que existam todos os $\inf$s e
$\sup$s. Mostre que
\begin{equation*}
\inf\, B \leq \inf\, A \leq \sup\, A \leq \sup\, B .
\end{equation*}
\end{exercise}

\begin{exercise}
Seja $S$ um conjunto ordenado. Seja $A \subset S$ e suponha que
$b$ seja uma cota superior de $A$. Suponha que $b \in A$. Mostre
que $b = \sup\, A$.
\end{exercise}

\begin{exercise}
Seja $S$ um conjunto ordenado. Seja $A \subset S$ não vazio e
limitado superiormente.
Suponha que $\sup\, A$ exista e
$\sup\, A \notin A$.
Mostre que $A$ contém um subconjunto infinito enumerável.
%In particular, $A$ is infinite.
\end{exercise}

\begin{exercise}
Encontre uma ordenação não usual do conjunto dos números naturais $\N$
tal que exista um subconjunto próprio não vazio $A \subsetneq \N$
e tal que $\sup\, A$ exista em $\N$, mas $\sup\, A \notin A$.
Para manter a notação clara, pode ser uma boa ideia usar uma notação diferente
para a ordenação não usual, como $n \prec m$.
\end{exercise}

\begin{exercise}
\pagebreak[2]
Seja $F \coloneqq \{ 0, 1, 2 \}$.
\begin{enumerate}[a)]
\item
Demonstre que existe uma única maneira de definir adição e multiplicação
para que $F$ seja um corpo, se $0$ e $1$ tiverem o significado usual dos
axiomas (A4) e (M4).
\item
Mostre que $F$ não pode ser um corpo ordenado.
\end{enumerate}
\end{exercise}

\begin{exercise} \label{exercise:dominatingb}
\pagebreak[2]
Seja $S$ um conjunto ordenado e
seja $A$ um subconjunto não vazio tal que $\sup \, A$ exista. Suponha que exista
$B \subset A$ tal que, sempre que $x \in A$, existe $y \in B$
com $x \leq y$. Mostre que $\sup \, B$ existe e $\sup \, B = \sup \, A$.
\end{exercise}

\begin{exercise}
Seja $D$ o conjunto ordenado de todas as palavras possíveis (não apenas palavras em inglês,
mas todas as sequências de letras de comprimento arbitrário)
formadas com o alfabeto latino, usando somente letras minúsculas. A ordem é a
ordem lexicográfica, como em um dicionário (por exemplo, \emph{aa} $<$ \emph{aaa} $<$ \emph{dog} $<$ \emph{door}).
Seja $A$ o subconjunto de $D$ que contém as palavras cuja
primeira letra é `a' (por exemplo, \emph{a} $\in A$, \emph{abcd} $\in A$).
Mostre que $A$ tem supremo e determine qual é.
\end{exercise}

\begin{exercise}
Sejam $F$ um corpo ordenado e $x,y,z,w \in F$.
\begin{enumerate}[a)]
\item
Demonstre a parte \ref{prop:bordfield:vii} de \propref{prop:bordfield}.
Isto é,
se $x \leq y$ e $z \leq w$, então $x+z \leq y+w$.
\item
Demonstre que,
se $x < y$ e $z \leq w$, então $x+z < y+w$.
\end{enumerate}
\end{exercise}

\begin{exercise}
Demonstre que todo corpo ordenado contém um subconjunto infinito enumerável.
\end{exercise}

\begin{exercise}
Seja $\N_{\infty} \coloneqq \N \cup \{ \infty \}$, onde os elementos de $\N$ são
ordenados entre si da maneira usual,
e $k < \infty$ para todo $k \in \N$. Mostre que $\N_{\infty}$ é um conjunto ordenado
e que todo subconjunto $E \subset \N_{\infty}$ tem supremo em
$\N_{\infty}$ (inclua também o caso do conjunto vazio).
Observação: o caso do conjunto vazio será diferente
se o conjunto ordenado for o dos números reais ou o dos números reais
estendidos, que aparecerão na próxima seção. É importante que, neste exercício,
o conjunto ordenado seja $\N_{\infty}$.
\end{exercise}

\begin{exercise}
Seja $S \coloneqq \{ a_k : k \in \N \} \cup \{ b_k : k \in \N \}$, ordenado de modo que
$a_k < b_j$ para todo $k$ e $j$,
$a_k < a_m$ sempre que $k < m$,
e $b_k > b_m$ sempre que $k < m$.
\begin{enumerate}[a)]
\item
Mostre que $S$ é um conjunto ordenado.
\item
Mostre que todo subconjunto de $S$ é limitado (superiormente e inferiormente).
\item
Encontre um subconjunto limitado de $S$ que não tenha menor cota superior.
\end{enumerate}
\end{exercise}

%%%%%%%%%%%%%%%%%%%%%%%%%%%%%%%%%%%%%%%%%%%%%%%%%%%%%%%%%%%%%%%%%%%%%%%%%%%%%%

\sectionnewpage

\section{O conjunto dos números reais} \label{sec:setofreals}

%mbxINTROSUBSECTION

\sectionnotes{2 aulas; os números reais estendidos são opcionais}

\subsection{O conjunto dos números reais}

Finalmente chegamos ao sistema dos números reais. Para simplificar, em vez de
construir o conjunto dos números reais a partir dos racionais, simplesmente
enunciamos sua existência como um teorema, sem demonstração. Observe que $\Q$ é um corpo
ordenado.

\begin{thm}
Existe um único\footnote{A unicidade é a menos de isomorfismo, mas queremos
evitar o uso excessivo de álgebra. Para nós, basta supor
que existe um conjunto dos números reais. Veja Rudin~\cite{Rudin:baby} para a
construção e mais detalhes.}
corpo ordenado $\R$ com a propriedade da
\hyperref[defn:lub]{menor cota superior}
tal que $\Q \subset \R$.
\end{thm}

Observe também que $\N \subset \Q$. Vimos que $1 > 0$.
Pela \hyperref[induction:thm]{indução} (exercício),
podemos provar que $n > 0$ para todo $n \in \N$.
Da mesma forma, verificamos afirmações simples sobre os números racionais.
Por exemplo, provamos que, se $n > 0$, então $\nicefrac{1}{n} > 0$.
Então $m < k$ implica $\nicefrac{m}{n} < \nicefrac{k}{n}$.

A análise consiste em demonstrar desigualdades, e
a proposição a seguir, ou alguma de suas muitas variações,
é a maneira como um analista demonstra uma desigualdade não estrita.

\begin{prop}
Se $x \in \R$ é tal que $x \leq \epsilon$ para todo
$\epsilon \in \R$ tal que
$\epsilon > 0$, então $x \leq 0$.
\end{prop}

\begin{proof}
Se $x > 0$, então $0 < \nicefrac{x}{2} < x$ (por quê?). Tome
$\epsilon = \nicefrac{x}{2}$ para obter uma contradição. Portanto, $x \leq 0$.
\end{proof}

Para $x$ não negativo, obtemos uma igualdade:
\emph{Se $x \geq 0$ é tal que $x \leq \epsilon$ para todo
$\epsilon > 0$, então $x = 0$.}
Uma versão comum usa o valor absoluto (veja \sectionref{sec:absval}):
\emph{Se $\sabs{x} \leq \epsilon$ para todo
$\epsilon > 0$, então $x = 0$.}
Para provar que $x \geq 0$, um analista pode demonstrar
que $x \geq -\epsilon$ para todo $\epsilon > 0$.
De agora em diante, quando dissermos $x \geq 0$ ou
$\epsilon > 0$, estaremos supondo automaticamente que $x \in \R$ e $\epsilon \in \R$.

A ideia por trás da proposição acima é que, sempre que tivermos dois números
reais $a < b$, existe outro
número real $c$ tal que
$a < c < b$. Existem infinitos valores de $c$ assim. Por exemplo, um deles é
$c = \frac{a+b}{2}$ (por quê?).
Usaremos esse fato no próximo exemplo.

A propriedade mais útil de $\R$ para os analistas
não é apenas ser um corpo ordenado, mas também ter a
\hyperref[defn:lub]{propriedade da menor cota superior}.
Essencialmente, queremos $\Q$, mas também queremos
tomar supremos (e ínfimos) à vontade.
Então partimos de $\Q$ e acrescentamos números suficientes para obter $\R$.

Já mencionamos que $\R$ contém elementos que não pertencem a $\Q$
por causa da \hyperref[defn:lub]{propriedade da menor cota superior}. Vamos
demonstrar isso.
Vimos que não existe raiz quadrada racional de dois. O conjunto
$\{ x \in \Q : x^2 < 2 \}$ implica a existência do número real
$\sqrt{2}$, embora esse fato exija algum trabalho. Veja também
\exerciseref{exercise:sqrt2QorR}.

\begin{example} \label{example:sqrt2}
Afirmação: \emph{Existe um único
$r \in \R$ positivo tal que $r^2 = 2$. Denotamos $r$ por $\sqrt{2}$.}

\begin{proof}
Considere o conjunto
$A \coloneqq \{ x \in \R : x^2 < 2 \}$. Primeiro mostraremos que $A$ é limitado superiormente e
não vazio. A desigualdade $x \geq 2$ implica $x^2 \geq 4$
(veja \exerciseref{exercise:squareineq}).
Portanto, se $x^2 < 2$, então $x < 2$. Logo, $A$ é limitado superiormente.
Como $1 \in A$, o conjunto $A$ não é vazio. Portanto, existe o supremo.

Seja $r \coloneqq \sup\, A$. Mostraremos que $r^2 = 2$ provando
que $r^2 \geq 2$ e $r^2 \leq 2$. É assim que os analistas demonstram
uma igualdade: provando duas desigualdades.
Já sabemos que $r \geq 1 > 0$.

%\medskip
A seguir, pode parecer que estamos tirando certas expressões
da cartola. Ao escrever uma demonstração como esta, naturalmente só
encontraríamos as expressões depois de explorar o que queremos provar. A
ordem em que escrevemos a demonstração não é necessariamente a ordem em que
a encontramos.

Vamos primeiro mostrar que $r^2 \geq 2$.
Tome um número positivo $s$ tal que $s^2 < 2$. Queremos encontrar um $h > 0$
tal que ${(s+h)}^2 < 2$.
Como $2-s^2 > 0$, temos $\frac{2-s^2}{2s+1} > 0$.
Escolha $h \in \R$ tal que
$0 < h < \frac{2-s^2}{2s+1}$.
Além disso, suponha que $h < 1$. Estimamos:
\begin{equation*}
\begin{aligned}
{(s+h)}^2 - s^2 & = h(2s + h) \\
 & < h(2s+1) & & \quad \bigl(\text{pois } h < 1\bigr) \\
 & < 2-s^2 & & \quad \bigl(\text{pois } h < \tfrac{2-s^2}{2s+1} \bigr) .
\end{aligned}
\end{equation*}
Portanto, ${(s+h)}^2 < 2$. Assim, $s+h \in A$, mas, como $h > 0$,
temos $s+h > s$. Logo, $s < r = \sup\, A$. Como $s$ era um número
positivo arbitrário tal que $s^2 < 2$, segue que $r^2 \geq 2$.

%\medskip

Agora tome um número positivo $s$ tal que
$s^2 > 2$. Queremos encontrar um $h > 0$ tal que
${(s-h)}^2 > 2$ e $s-h$ ainda seja positivo.
Como $s^2-2 > 0$, temos $\frac{s^2-2}{2s} > 0$.
Defina $h \coloneqq \frac{s^2-2}{2s}$,
e verifique que $s-h=s-\frac{s^2-2}{2s} = \frac{s}{2}+\frac{1}{s} > 0$.
Estimamos:
\begin{equation*}
\begin{aligned}
s^2 - {(s-h)}^2 & = 2sh - h^2 \\
 & < 2sh & & \quad \bigl( \text{pois } h^2 > 0 \text{ as } h \neq 0 \bigr) \\
 & = s^2-2 & & \quad \bigl( \text{pois } h = \tfrac{s^2-2}{2s} \bigr) .
\end{aligned}
\end{equation*}
Subtraindo $s^2$ de ambos os lados e multiplicando por $-1$, obtemos
${(s-h)}^2 > 2$. Portanto, $s-h \notin A$.

Além disso, se $x \geq s-h$,
então $x^2 \geq {(s-h)}^2 > 2$ (pois $x > 0$ e $s-h > 0$), e assim $x \notin A$.
Logo,
$s-h$ é uma cota superior de $A$. No entanto, $s-h < s$, ou, em outras
palavras, $s > r = \sup\, A$. Portanto, $r^2 \leq 2$.

\medskip

Juntas, $r^2 \geq 2$ e $r^2 \leq 2$ implicam
$r^2 = 2$. A parte da existência está concluída. Ainda precisamos
provar a unicidade. Suponha que $s \in \R$ seja tal que $s^2 = 2$ e $s > 0$.
Assim, $s^2 = r^2$. No entanto, se $0 < s < r$, então $s^2 < r^2$. Da mesma forma,
$0 < r < s$ implica $r^2 < s^2$. Portanto, $s = r$.
\end{proof}
\end{example}

O número $\sqrt{2} \notin \Q$. O conjunto
$\R \setminus \Q$ é chamado de conjunto dos números
\emph{\myindex{irracionais}}. Acabamos de provar que $\R \setminus \Q$
não é vazio. E, como veremos, ele não é apenas não vazio: é muito
grande mesmo.

Usando a mesma técnica acima, podemos mostrar que existe um número real positivo
$x^{1/n}$ para todo $n\in \N$ e todo $x > 0$.
Isto é, para cada $x > 0$,
existe um único número real positivo $r$ tal que $r^n = x$.
A demonstração fica como exercício.

\subsection{Propriedade arquimediana}

Como vimos, há muitos números reais em qualquer intervalo. Mas também
há infinitos números racionais em qualquer intervalo. O resultado a seguir é
um dos fatos fundamentais sobre os números reais.
As duas partes do próximo teorema são, na verdade, equivalentes, embora isso
talvez não pareça evidente à primeira vista.


\begin{thm} \label{thm:arch}
%FIXMEevillayouthack
\pagebreak[3]
\leavevmode
\begin{enumerate}[(i)]
\item \label{thm:arch:i} \emph{(\myindex{propriedade arquimediana})}%
\footnote{Em homenagem ao matemático da Grécia Antiga
\href{https://en.wikipedia.org/wiki/Archimedes}{Arquimedes de Siracusa}
(c.\ 287 a.C. -- c.\ 212 a.C.)\@.
Esta propriedade é o Axioma V da obra de Arquimedes
\myquote{Sobre a esfera e o cilindro} (225 a.C.)\@.}
Se $x, y \in \R$ e
$x > 0$, então existe $n \in \N$ tal que
\begin{equation*}
nx > y .
\end{equation*}
\item \label{thm:arch:ii} \emph{($\Q$ é denso em $\R$\index{densidade dos números racionais})} Se $x, y \in \R$ e
$x < y$, então existe $r \in \Q$ tal que
$x < r < y$.
\end{enumerate}
\end{thm}

\begin{proof}
Vamos provar \ref{thm:arch:i}. Divida tudo por $x$.
Então \ref{thm:arch:i} afirma que, para todo número real $t\coloneqq \nicefrac{y}{x}$,
podemos encontrar $n \in \N$ tal que $n > t$. Em outras palavras,
\ref{thm:arch:i} afirma que $\N \subset \R$ não é limitado superiormente.
Suponha, por absurdo, que $\N$ seja limitado superiormente. Seja $b \coloneqq \sup \N$.
O número $b-1$ não pode ser uma cota superior de $\N$, pois é estritamente
menor que $b$ (a menor cota superior). Portanto, existe $m \in \N$ tal que $m > b-1$.
Some 1 aos dois lados para obter $m+1 > b$, uma contradição, pois $b$ é uma
cota superior.

\begin{myfigureht}
\myincludepdft{figdensofQ}{%
Um diagrama da reta real horizontal com vários traços verticais altos
igualmente espaçados, cuja distância entre eles é 1 sobre n. O segundo traço
a partir da esquerda está marcado como m menos 1, tudo sobre n.
Depois desse traço, aparece o número x, marcado com um traço curto.
O próximo traço alto está marcado como m sobre n.
Em seguida aparece o número y e depois um traço alto marcado como
m mais 1, tudo sobre n. Mais dois traços altos aparecem depois.}
\caption{Ideia da demonstração da densidade de $\Q$: encontre $n$ tal que $y-x >
\nicefrac{1}{n}$ e, em seguida, tome o menor $m$ tal que $\nicefrac{m}{n} > x$.\label{figdensofQ}}
\end{myfigureht}
Vamos demonstrar \ref{thm:arch:ii}.
Veja a \figureref{figdensofQ}
para visualizar a ideia da demonstração.
Suponha primeiro que $x \geq 0$.
Observe que $y-x > 0$.
Por \ref{thm:arch:i}, existe $n \in \N$ tal que
\begin{equation*}
n(y-x) > 1
\qquad \text{ou} \qquad
y-x > \nicefrac{1}{n}.
\end{equation*}
Novamente por \ref{thm:arch:i}, o conjunto
$A \coloneqq \{ k \in \N : k > nx \}$ não é vazio. Pela
\hyperlink{wop:link}{propriedade da boa ordenação}
de $\N$, $A$ tem um menor elemento $m$.
Como $m \in A$, temos $m > nx$.
Divida tudo por $n$ para obter $x < \nicefrac{m}{n}$.
Como $m$ é o menor
elemento de $A$, $m-1 \notin A$.
Se $m > 1$, então $m-1 \in \N$, mas $m-1 \notin A$ e, portanto, $m-1 \leq nx$.
Se $m=1$,
então $m-1 = 0$, e ainda temos $m-1 \leq nx$, pois $x \geq 0$.
Em outras palavras,
\begin{equation*}
m-1 \leq nx \qquad \text{ou} \qquad m \leq nx+1 .
\end{equation*}
Por outro lado,
de $n(y-x) > 1$ obtemos $ny > 1+nx$.
Logo, $ny > 1+nx \geq m$, e, portanto, $y > \nicefrac{m}{n}$.
Juntando tudo, obtemos $x < \nicefrac{m}{n} < y$.
Assim, tome $r = \nicefrac{m}{n}$.

Agora suponha que $x < 0$. Se $y > 0$, basta tomar $r=0$. Se
$y \leq 0$, então $0 \leq -y < -x$, e encontramos
um número racional $q$ tal que $-y < q < -x$. Então, tome $r = -q$.
\end{proof}

Enunciemos e provemos um corolário simples, mas útil, da
\hyperref[thm:arch:i]{propriedade arquimediana}.

\begin{cor}
$\inf \{ \nicefrac{1}{n} : n \in \N \} = 0$.
Veja a \figureref{fig:oneovernset}.
\end{cor}

\begin{proof}
Seja $A \coloneqq \{ \nicefrac{1}{n} : n \in \N \}$. É evidente que $A$ não é vazio.
Além disso,
$\nicefrac{1}{n} > 0$ para todo $n \in \N$, então $0$ é uma cota inferior
e $b \coloneqq \inf\, A$ existe.
Como $0$ é uma cota inferior, $b \geq 0$.
Tome um $a > 0$ arbitrário. Pela
\hyperref[thm:arch:i]{propriedade arquimediana}, existe $n$ tal que
$na > 1$, isto é, $a > \nicefrac{1}{n} \in A$. Portanto,
$a$ não pode ser uma cota inferior de $A$. Logo, $b=0$.
\end{proof}

\begin{myfigureht}
\myincludegraphics{oneovernset}{%
Um diagrama da reta real horizontal com traços verticais indicando
os números 1, 1 sobre 2, 1 sobre 3, 1 sobre 4 etc. Os traços ficam
cada vez mais próximos de 0, que também está marcado. Ao chegar a 0,
eles são tão espessos que só vemos ali um retângulo preto.}
\caption{O conjunto $\{ \nicefrac{1}{n} : n \in \N \}$ e seu ínfimo $0$.\label{fig:oneovernset}}
\end{myfigureht}



\subsection{Usando supremo e ínfimo}

Supremos e ínfimos são
compatíveis com as operações algébricas. Para um conjunto $A \subset \R$ e
$x \in \R$, definimos
\begin{align*}
x + A & \coloneqq \{ x+y \in \R : y \in A \} , \\
xA & \coloneqq \{ xy \in \R : y \in A \} .
\end{align*}
Por exemplo, se $A = \{ 1,2,3 \}$, então $5+A = \{ 6,7,8 \}$ e $3A = \{ 3,6,9
\}$.

\begin{prop} \label{prop:supinfalg}
Seja $A \subset \R$ não vazio.
\begin{enumerate}[(i)]
\item Se $x \in \R$ e $A$ é limitado superiormente, então $\sup (x+A) = x + \sup\, A$.
\item Se $x \in \R$ e $A$ é limitado inferiormente, então $\inf (x+A) = x + \inf\, A$.
\item Se $x > 0$ e $A$ é limitado superiormente, então $\sup (xA) = x ( \sup\, A )$.
\item Se $x > 0$ e $A$ é limitado inferiormente, então $\inf (xA) = x ( \inf\, A )$.
\item Se $x < 0$ e $A$ é limitado inferiormente, então $\sup (xA) = x ( \inf\, A )$.
\item Se $x < 0$ e $A$ é limitado superiormente, então $\inf (xA) = x ( \sup\, A )$.
\end{enumerate}
\end{prop}

Observe que a multiplicação de um conjunto por um número negativo troca supremos e ínfimos.
Além disso, como a proposição implica que
o supremo (respectivamente, o ínfimo) de $x+A$ ou $xA$ existe,
ela também implica que $x+A$ ou $xA$
é não vazio e limitado superiormente (respectivamente, inferiormente).

\begin{proof}
Vamos provar apenas a primeira afirmação. As demais ficam como exercícios.

Suponha que $b$ seja uma cota superior de $A$. Isto é, $y \leq b$ para todo $y \in A$.
Então $x+y \leq x+b$ para todo $y \in A$, de modo que $x+b$ é uma cota
superior de $x+A$. Em particular, se $b = \sup\, A$, então
\begin{equation*}
\sup (x+A) \leq x+b = x+ \sup\, A .
\end{equation*}

A desigualdade oposta se prova de maneira semelhante. Se $c$ é uma cota superior de $x+A$,
então $x+y \leq c$
para todo $y \in A$ e, portanto, $y \leq c-x$ para todo $y \in A$.
Logo, $c-x$ é uma cota superior de $A$.
Se $c = \sup (x+A)$, então
\begin{equation*}
\sup\, A \leq c-x = \sup (x+A) -x .
\end{equation*}
A conclusão segue.
\end{proof}

Às vezes precisamos aplicar supremo ou ínfimo duas vezes. Veja este exemplo.

\begin{prop} \label{infsupineq:prop}
Sejam $A, B \subset \R$ conjuntos não vazios tais que $x \leq y$ sempre que $x \in A$ e
$y \in B$. Então $A$ é limitado superiormente, $B$ é limitado inferiormente e $\sup\, A \leq \inf\, B$.
\end{prop}

\begin{proof}
Todo $x \in A$ é uma cota inferior de $B$.
Portanto, $x \leq \inf\, B$ para todo $x \in A$,
de modo que $\inf\, B$ é uma cota superior de $A$.
Logo, $\sup\, A \leq \inf\, B$.
\end{proof}

É preciso ter cuidado com desigualdades estritas e com a aplicação de supremos e ínfimos.
Observe que $x < y$ sempre que $x \in A$ e
$y \in B$ ainda implica apenas $\sup\, A \leq \inf\, B$, e não uma desigualdade
estrita.
Por exemplo, tome $A \coloneqq \{ 0 \}$ e $B \coloneqq \{ \nicefrac{1}{n}
: n \in \N \}$.
Então $0 < \nicefrac{1}{n}$
para todo $n \in \N$. No entanto, $\sup\, A = 0$ e $\inf\, B = 0$.
Esse importante detalhe sutil aparece com frequência.

\medskip
\pagebreak[1]

Deixamos ao leitor a demonstração do seguinte
resultado, que usamos com frequência.
Um resultado semelhante vale para ínfimos.

\begin{prop} \label{prop:existsxepsfromsup}
Se $S \subset \R$ é não vazio e limitado superiormente,
então, para todo $\epsilon > 0$, existe $x \in S$ tal
que $(\sup\, S) - \epsilon < x \leq \sup\, S$.
\end{prop}


Para facilitar ainda mais o uso de supremos e ínfimos, talvez queiramos
escrever $\sup\, A$ e $\inf\, A$ sem nos preocupar se $A$ é
limitado e não vazio. Vamos fazer as seguintes definições naturais.

\begin{defn}
Seja $A \subset \R$ um conjunto.
\begin{enumerate}[(i)]
\item Se $A$ é vazio, então $\sup\, A \coloneqq -\infty$.
\item Se $A$ não é limitado superiormente, então $\sup\, A \coloneqq \infty$.
\item Se $A$ é vazio, então $\inf\, A \coloneqq \infty$.
\item Se $A$ não é limitado inferiormente, então $\inf\, A \coloneqq -\infty$.
\end{enumerate}
\end{defn}

As afirmações da definição acima
decorrem naturalmente se, em vez de trabalhar no
conjunto ordenado $\R$, tomarmos supremos e ínfimos
no conjunto ordenado
$\R^* \coloneqq \R \cup \{ -\infty , \infty\}$
\glsadd{not:extreal},\glsadd{not:infinity},
onde definimos
\begin{equation*}
-\infty < \infty \quad \text{e} \quad
-\infty < x \quad \text{e} \quad
x < \infty \quad \text{para todo $x \in \R$}.
\end{equation*}
O conjunto $\R^*$ é chamado de conjunto dos \emph{\myindex{números reais estendidos}}.
Por conveniência, às vezes tratamos $\infty$ e $-\infty$ como se fossem
números, mas não permitimos operações aritméticas arbitrárias com eles.
É possível definir algumas operações aritméticas em $\R^*$. A maioria delas
se estende de maneira evidente, mas precisamos deixar
$\infty-\infty$, $0 \cdot (\pm\infty)$ e $\frac{\pm\infty}{\pm\infty}$
sem definição.
Evitamos
usar essas operações;
elas levam a erros fáceis, pois $\R^*$ não é um corpo.
Agora podemos tomar supremos e ínfimos sem nos preocupar com conjuntos vazios ou
não limitados. Neste livro, evitamos em grande parte
usar $\R^*$ fora dos exercícios e deixamos tais generalizações para o leitor interessado.

\subsection{Máximos e mínimos}

Pelo \exerciseref{exercise:finitesethasminmax},
todo conjunto finito de números tem supremo e ínfimo,
e ambos pertencem ao próprio conjunto.
Neste caso, em geral, não usamos as palavras supremo ou ínfimo.
Quando um conjunto $A$ de números reais é limitado superiormente e
$\sup\, A \in A$, podemos usar a palavra
\emph{\myindex{máximo}} e a notação $\max\, A$\glsadd{not:max} para indicar o supremo.
De maneira semelhante, para o ínfimo: quando $A$ é limitado inferiormente
e $\inf\, A \in A$, podemos usar a
palavra \emph{\myindex{mínimo}} e a notação $\min\, A$.\glsadd{not:min}
Por exemplo,
\begin{align*}
& \max \{ 1,2.4,\pi,100 \} = 100 , \\
& \min \{ 1,2.4,\pi,100 \} = 1 .
\end{align*}
Embora escrever $\sup$ e $\inf$ seja tecnicamente
correto nesta situação, em geral se usam $\max$ e
$\min$ para enfatizar que o supremo ou o ínfimo
pertence ao próprio conjunto, especialmente quando o conjunto é finito.

\subsection{Exercícios}

\begin{exercise}
Demonstre que,
se $t > 0$ ($t \in \R$), então existe $n \in \N$ tal que
$\dfrac{1}{n^2} < t$.
\end{exercise}

\begin{exercise}
Demonstre que,
se $t \geq 0$ ($t \in \R$), então existe $n \in \N$ tal que $n-1 \leq t < n$.
\end{exercise}

\begin{exercise}
Complete a demonstração de \propref{prop:supinfalg}.
\end{exercise}

\begin{exercise}
Sejam $x, y \in \R$. Suponha que $x^2 + y^2 = 0$. Demonstre que $x = 0$ e $y = 0$.
\end{exercise}

\begin{exercise}
Mostre que $\sqrt{3}$ é irracional.
\end{exercise}

\begin{exercise}
Seja $n \in \N$.
Mostre que $\sqrt{n}$ ou é um inteiro ou é
irracional.
\end{exercise}

\begin{exercise}
Demonstre a \emph{\myindex{desigualdade entre as médias aritmética e geométrica}}.
Para dois números reais positivos $x,y$,
\begin{equation*}
\sqrt{xy} \leq \frac{x+y}{2} .
\end{equation*}
Além disso, há igualdade se, e somente se, $x=y$.
\end{exercise}

\begin{exercise}
Mostre que, para todo par de números reais $x$ e $y$ tais que $x < y$, existe
um número irracional $s$ tal que $x < s < y$. Dica:
aplique a densidade de $\Q$ a $\dfrac{x}{\sqrt{2}}$ e
$\dfrac{y}{\sqrt{2}}$, mas evite o zero.
\end{exercise}

\begin{exercise} \label{exercise:supofsum}
Sejam $A$ e $B$ dois conjuntos não vazios e limitados de números reais. Seja
$C \coloneqq \{ a+b : a \in A, b \in B \}$.
Mostre que $C$ é limitado e que
\begin{equation*}
\sup\,C = \sup\,A + \sup\,B
\qquad \text{e} \qquad
\inf\,C = \inf\,A + \inf\,B .
\end{equation*}
\end{exercise}

\begin{exercise}
Sejam $A$ e $B$ dois conjuntos não vazios e limitados de números reais não negativos.
Defina o conjunto
$C \coloneqq \{ ab : a \in A, b \in B \}$.
Mostre que $C$ é limitado e que
\begin{equation*}
\sup\,C = (\sup\,A )( \sup\,B)
\qquad \text{e} \qquad
\inf\,C = (\inf\,A )( \inf\,B).
\end{equation*}
\end{exercise}

\begin{exercise}[Difícil] \label{exercise:rootexistshard}
Dados $x > 0$ e $n \in \N$, mostre que existe um único número real positivo
$r$ tal que $x = r^n$. Em geral, denotamos $r$ por $x^{1/n}$.
\end{exercise}

\begin{exercise}[Fácil]
Demonstre \propref{prop:existsxepsfromsup}.
\end{exercise}

\begin{exercise} \label{exercise:bernoulliineq}
Demonstre a \emph{\myindex{desigualdade de Bernoulli}}%
\footnote{%
Em homenagem ao matemático suíço
\href{https://en.wikipedia.org/wiki/Jacob_Bernoulli}{Jacob Bernoulli}
(1655--1705).}%
: se $1+x > 0$, então,
para todo $n \in \N$, temos $(1+x)^n \geq 1+nx$.
\end{exercise}

\begin{exercise} \label{exercise:sqrt2QorR}
Demonstre que $\sup \{ x \in \Q : x^2 < 2 \} = \sup \{ x \in \R : x^2 < 2 \}$.
\end{exercise}

\begin{exercise} \label{exercise:Dedekind}
\leavevmode
\begin{enumerate}[a)]
\item
Demonstre que, dado $y \in \R$, temos $\sup \{ x \in \Q : x < y \} = y$.
\item
Seja $A \subset \Q$ um conjunto não vazio e limitado superiormente tal que, sempre que $x
\in A$ e $t \in \Q$ com $t < x$, então $t \in A$. Suponha também
que $\sup\, A \notin A$. Mostre que existe $y \in \R$ tal que
$A = \{ x \in \Q : x < y \}$. Um conjunto como $A$ é chamado de
\emph{\myindex{corte de Dedekind}}.
\item
Mostre que existe uma bijeção entre $\R$ e os cortes de Dedekind.
\end{enumerate}
Observação: Dedekind usou conjuntos como os da parte b) em sua construção dos
números reais.
\end{exercise}

\begin{exercise}
Demonstre que, se $A \subset \Z$ é um subconjunto não vazio e limitado inferiormente, então
$A$ tem um menor elemento.
Agora descreva por que essa afirmação simplificaria a demonstração do
\thmref{thm:arch}, parte \ref{thm:arch:ii}, de modo que não seria preciso supor
$x \geq 0$.
\end{exercise}

\begin{samepage}
\begin{exercise} \label{exercise:realpower}
Suponha que saibamos que $x^{1/n}$ existe para todo $x > 0$ e todo $n \in
\N$ (veja acima \exerciseref{exercise:rootexistshard}).
Para inteiros $p$ e $q > 0$ tais que $\nicefrac{p}{q}$ esteja em forma irredutível,
defina $x^{p/q} \coloneqq
{(x^{1/q})}^p$.
\begin{enumerate}[a)]
\item
Mostre que a potência está bem definida mesmo
se a fração não estiver em forma irredutível: se $\nicefrac{p}{q} =
\nicefrac{m}{k}$, onde $m$ e $k > 0$ são inteiros, então
${(x^{1/q})}^p = {(x^{1/k})}^m$.
\item
Sejam $x$ e $y$ dois números positivos e $r$ um número racional.
Supondo que $r > 0$, mostre que
$x < y$ se, e somente se, $x^r < y^r$. Em seguida, suponha que $r < 0$ e mostre:
$x < y$ se, e somente se, $x^r > y^r$.
\item
Suponha que $x > 1$ e que $r,s$ sejam números racionais tais que $r < s$.
Mostre que $x^r < x^s$. Se $0 < x < 1$ e $r < s$, mostre que $x^r > x^s$.
Dica: escreva $r$ e $s$ com o mesmo denominador.
\item
(Desafio)%
\footnote{Em
\sectionref{sec:logandexp}
definiremos a exponencial e o logaritmo, e
definiremos $x^z \coloneqq \exp(z \ln x)$. Assim, teremos
ferramentas suficientes para demonstrar essas afirmações com muito mais facilidade. Neste
ponto, porém, ainda não dispomos dessas ferramentas.}
Para um número irracional $z \in \R \setminus \Q$ e $x > 1$, defina
$x^z \coloneqq \sup \{ x^r : r \leq z, r \in \Q \}$;
para $x=1$, defina $1^z = 1$;
e, para $0 < x < 1$, defina
$x^z \coloneqq \inf \{ x^r : r \leq z, r \in \Q \}$.
Demonstre as duas afirmações da parte b) para todo número real $z$.
\end{enumerate}
\end{exercise}
\end{samepage}



%%%%%%%%%%%%%%%%%%%%%%%%%%%%%%%%%%%%%%%%%%%%%%%%%%%%%%%%%%%%%%%%%%%%%%%%%%%%%%

\sectionnewpage

\section{Valor absoluto e funções limitadas} \label{sec:absval}

\sectionnotes{0,5--1 aula}

Um conceito que encontraremos repetidamente é o
\emph{\myindex{valor absoluto}}.
Pense no valor absoluto como o \myquote{tamanho} de um número real.
Eis a definição formal:
\glsadd{not:abs}
\begin{equation*}
\sabs{x} \coloneqq
\begin{cases}
x & \text{se } x \geq 0, \\
-x & \text{se } x < 0 .
\end{cases}
\end{equation*}
A proposição a seguir apresenta
as principais propriedades do valor
absoluto.

\begin{prop} \label{prop:absbas}
\leavevmode
\begin{enumerate}[(i)]
\item \label{prop:absbas:i} $\sabs{x} \geq 0$, e $\sabs{x}=0$ se, e somente se, $x = 0$.
\item \label{prop:absbas:ii} $\sabs{-x} = \sabs{x}$ para todo $x \in \R$.
\item \label{prop:absbas:iii} $\sabs{xy} = \sabs{x}\sabs{y}$ para todos $x,y \in \R$.
\item \label{prop:absbas:iv} $\sabs{x}^2 = x^2$ para todo $x \in \R$.
\item \label{prop:absbas:v} $\sabs{x} \leq y$ se, e somente se, $-y \leq x \leq y$.
\item \label{prop:absbas:vi} $-\sabs{x} \leq x \leq \sabs{x}$ para todo $x \in \R$.
\end{enumerate}
\end{prop}

\begin{proof}
\ref{prop:absbas:i}:
Suponha primeiro que $x \geq 0$.
Então $\sabs{x} = x \geq 0$. Além disso, $\sabs{x} = x = 0$ se, e somente
se, $x=0$.
Por outro lado, se $x < 0$, então $\sabs{x} = -x > 0$, e $\sabs{x}$ nunca é zero.

\medskip

\ref{prop:absbas:ii}: Se $x > 0$, então $-x < 0$ e, portanto, $\sabs{-x} = -(-x) = x =
\sabs{x}$. O mesmo vale quando $x < 0$ ou $x = 0$.

\medskip

\ref{prop:absbas:iii}:
Se $x$ ou $y$ é zero, o resultado é imediato. Se $x$ e
$y$ são ambos positivos, então $\sabs{x}\sabs{y} = xy$. O produto de números
positivos é positivo, logo $xy$ é
positivo e $xy = \sabs{xy}$.
Se $x$ e $y$ são ambos negativos,
então $xy=(-x)(-y)$ é positivo pelo mesmo motivo e, novamente, $\sabs{xy}=xy$. Além disso,
$\sabs{x}\sabs{y} = (-x)(-y) = xy$.
Logo, o resultado vale quando $x$ e $y$ têm o mesmo sinal.
Se $x > 0$ e $y < 0$, então $\sabs{x}\sabs{y} = x(-y) = -(xy)$. Agora
$xy$ é negativo e $\sabs{xy} = -(xy)$, e o resultado vale. O mesmo ocorre quando
$x < 0$ e $y > 0$.

\medskip

\ref{prop:absbas:iv}:
É imediato se $x \geq 0$. Se $x < 0$, então $\sabs{x}^2 = {(-x)}^2 =
x^2$.

\medskip

\ref{prop:absbas:v}:
Suponha que $\sabs{x} \leq y$.
Se $x \geq 0$, então $x \leq y$.
Segue que $y \geq 0$, e portanto $-y \leq 0 \leq x$.
Assim, vale $-y \leq x \leq y$.
Se $x < 0$, então $\sabs{x} \leq y$ significa que $-x \leq y$.
Multiplicando ambos os lados por $-1$, obtemos $x \geq -y$.
Novamente, $y \geq 0$ e, portanto, $y \geq 0 > x$.
Logo, $-y \leq x \leq y$.

Por outro lado, suponha que $-y \leq x \leq y$.
Se $x \geq 0$, então $x \leq y$ equivale a $\sabs{x} \leq y$.
Se $x < 0$, então $-y \leq x$ implica $(-x) \leq y$,
o que equivale a $\sabs{x} \leq y$.

\medskip

\ref{prop:absbas:vi}:
Aplique \ref{prop:absbas:v} com $y = \sabs{x}$.
\end{proof}

Uma propriedade usada com frequência suficiente para merecer um nome é a
\emph{\myindex{desigualdade triangular}}.

\begin{prop}[Desigualdade Triangular]
$\sabs{x+y} \leq \sabs{x}+\sabs{y}$
para todos $x, y \in \R$.
\end{prop}

\begin{proof}
A \propref{prop:absbas} dá
$- \sabs{x} \leq x \leq \sabs{x}$ e
$- \sabs{y} \leq y \leq \sabs{y}$. Somando essas duas desigualdades, obtemos
\begin{equation*}
- \bigl(\sabs{x}+\sabs{y}\bigr) \leq x+y \leq \sabs{x}+ \sabs{y} .
\end{equation*}
Aplique \propref{prop:absbas} novamente para concluir que
$\sabs{x+y} \leq \sabs{x}+\sabs{y}$.
\end{proof}

Há outras versões da desigualdade triangular usadas com frequência.

\begin{cor}
Sejam $x,y \in \R$.
\begin{enumerate}[(i)]
\item \emph{(\myindex{desigualdade triangular reversa})}
~
$\babs{ \, \sabs{x}-\sabs{y} \, } \leq \sabs{x-y}$.
\item $\sabs{x-y} \leq \sabs{x}+\sabs{y}$.
\end{enumerate}
\end{cor}

\begin{proof}
Substituamos $x=a-b$ e $y=b$ na desigualdade triangular usual para obter
\begin{equation*}
\sabs{a} = \sabs{a-b+b} \leq \sabs{a-b} + \sabs{b} ,
\end{equation*}
ou $\sabs{a}-\sabs{b} \leq \sabs{a-b}$. Trocando os papéis de $a$ e $b$,
obtemos
$\sabs{b}-\sabs{a} \leq \sabs{b-a} = \sabs{a-b}$. Aplicando
\propref{prop:absbas}, obtemos a desigualdade triangular reversa.

A segunda afirmação do corolário decorre da desigualdade triangular
usual, substituindo $y$ por $-y$ e observando que $\sabs{-y} =
\sabs{y}$.
\end{proof}

\begin{cor}
Sejam $x_1, x_2, \ldots, x_n \in \R$. Então
\begin{equation*}
\sabs{x_1 + x_2 + \cdots + x_n} \leq
\sabs{x_1} + \sabs{x_2} + \cdots + \sabs{x_n} .
\end{equation*}
\end{cor}

\begin{proof}
Procedemos por \hyperref[induction:thm]{indução}.
A conclusão vale trivialmente para $n=1$ e,
para $n=2$, é a desigualdade triangular usual. Suponha que o corolário
valha para $n$. Considere $n+1$ números $x_1,x_2,\ldots,x_{n+1}$.
Primeiro aplique a desigualdade triangular usual e, em seguida, a hipótese
de indução:
\begin{equation*}
\begin{split}
\sabs{x_1 + x_2 + \cdots + x_n + x_{n+1}} & \leq
\sabs{x_1 + x_2 + \cdots + x_n} + \sabs{x_{n+1}} \\
& \leq
\sabs{x_1} + \sabs{x_2} + \cdots + \sabs{x_n} + \sabs{x_{n+1}} .  \qedhere
\end{split}
\end{equation*}
\end{proof}

Vejamos um exemplo do uso da desigualdade triangular.

\begin{example}
Encontre um número $M$ tal que $\sabs{x^2-9x+1} \leq M$ para todo $-1 \leq x \leq
5$.

Usando a desigualdade triangular, escrevemos
\begin{equation*}
\sabs{x^2-9x+1} \leq \sabs{x^2}+\sabs{-9x}+\sabs{1}
=
\sabs{x}^2+9\sabs{x}+1 .
\end{equation*}
A expressão
$\sabs{x}^2+9\sabs{x}+1$ é máxima quando $\sabs{x}$ é máximo (por quê?). No intervalo
dado, $\sabs{x}$ é máximo quando $x=5$, e então $\sabs{x}=5$. Uma
possibilidade para $M$ é
\begin{equation*}
M = 5^2+9(5)+1 = 71 .
\end{equation*}
Há, naturalmente, outros valores de $M$ que servem. A cota 71
é muito maior do que
precisa ser, mas não pedimos o melhor valor possível de $M$, apenas um que funcione.
\end{example}

O último exemplo nos leva ao conceito de funções limitadas.

\begin{defn}
Suponha que $f \colon D \to \R$ seja uma função. Dizemos que $f$ é
\emph{limitada}\index{função limitada}
se existe um número $M$
tal que $\babs{f(x)} \leq M$ para todo $x \in D$.
\end{defn}

No exemplo, provamos que $x^2-9x+1$ é limitada quando considerada como uma
função em $D = \{ x : -1 \leq x \leq 5 \}$. Por outro lado,
se considerarmos o mesmo polinômio como uma função em toda a reta real $\R$,
ele não é limitado.

\begin{myfigureht}
\myincludegraphics{boundedfunc}{%
Um diagrama do gráfico de uma função.
O domínio D está marcado no eixo horizontal, e uma curva senoidal
representando o gráfico de uma função nesse domínio é desenhada. O valor
máximo da função está indicado por uma linha horizontal tracejada que toca
o gráfico e está marcada como supremo de f em D. Da mesma forma, o
mínimo está marcado como ínfimo de f em D. Duas outras linhas tracejadas
pontilhadas estão marcadas com M e menos M. A linha marcada M está acima do
gráfico e a linha marcada menos M está abaixo. No eixo vertical, está marcado
o conjunto dos valores que a função assume, f de D.}
\caption{Exemplo de uma função limitada, uma cota $M$ e seu supremo e ínfimo.\label{boundedfuncfig}}
\end{myfigureht}
Para uma função $f \colon D \to \R$, escrevemos
(veja a \figureref{boundedfuncfig} para
um exemplo)
\glsadd{not:supf}\glsadd{not:inff}
\begin{equation*}
\sup_{x \in D} f(x) \coloneqq \sup\, f(D) \qquad \text{e} \qquad
\inf_{x \in D} f(x) \coloneqq \inf\, f(D) .
\end{equation*}
Às vezes, também substituímos o \myquote{$x \in D$} por uma expressão.
Por exemplo, como antes, se $f(x) = x^2-9x+1$, para $-1 \leq x \leq 5$,
um pouco de cálculo mostra que
\begin{equation*}
\sup_{x \in D} f(x) =
\sup_{-1 \leq x \leq 5} ( x^2 -9x+1 ) = 11,
\qquad
\inf_{x \in D} f(x) =
\inf_{-1 \leq x \leq 5} ( x^2 -9x+1 ) = \nicefrac{-77}{4} .
\end{equation*}



\begin{prop} \label{prop:funcsupinf}
Se $f \colon D \to \R$ e $g \colon D \to \R$ ($D$ não vazio) são
funções limitadas\footnote{A hipótese de limitação é adotada por simplicidade;
ela pode ser removida se usarmos os números reais estendidos.}
e
\begin{equation*}
f(x) \leq g(x) \qquad \text{para todo } x \in D,
\end{equation*}
então
\begin{equation} \label{prop:funcsupinf:eq}
\sup_{x \in D} f(x) \leq \sup_{x \in D} g(x)
\qquad \text{e} \qquad
\inf_{x \in D} f(x) \leq \inf_{x \in D} g(x) .
\end{equation}
\end{prop}

Tenha cuidado com as variáveis. O $x$ no lado esquerdo da
desigualdade em \eqref{prop:funcsupinf:eq}
é diferente do $x$ no lado direito. Na verdade,
deveríamos pensar, por exemplo, na primeira desigualdade como
\begin{equation*}
\sup_{x \in D} f(x) \leq \sup_{y \in D} g(y) .
\end{equation*}
Vamos provar essa desigualdade. Se $b$ é uma cota superior de $g(D)$, então
$f(x) \leq g(x) \leq b$ para todo $x \in D$; logo, $b$ também é uma cota superior
de $f(D)$,
ou seja, $f(x) \leq b$ para todo $x \in D$.
Tome a menor cota superior de $g(D)$ para concluir que, para todo $x \in D$,
\begin{equation*}
f(x) \leq \sup_{y \in D} g(y) .
\end{equation*}
Portanto,
$\sup_{y \in D} g(y)$ é uma cota superior de $f(D)$ e, assim, maior ou
igual à menor cota superior de $f(D)$.
\begin{equation*}
\sup_{x \in D} f(x) \leq \sup_{y \in D} g(y) .
\end{equation*}
A segunda desigualdade (a afirmação sobre o ínfimo) fica como exercício
(\exerciseref{exercise:finishfuncsupinf}).

\medskip

Um erro comum é concluir que
\begin{equation} \label{rn:av:ltnottrue}
\sup_{x \in D} f(x) \leq \inf_{y \in D} g(y) .
\end{equation}
A desigualdade \eqref{rn:av:ltnottrue} não decorre da hipótese
da proposição acima. Para obter essa desigualdade mais forte,
precisamos da hipótese mais forte
\begin{equation*}
f(x) \leq g(y) \qquad \text{para todo } x \in D \text{ and } y \in D.
\end{equation*}
A demonstração, assim como um contraexemplo, fica como exercício
(\exerciseref{exercise:supinffuncineq}).

\subsection{Exercícios}

\begin{exercise}
Mostre que
$\sabs{x-y} < \epsilon$ se, e somente se, $x-\epsilon < y < x+\epsilon$.
\end{exercise}

\begin{exercise}
Mostre: \qquad
%\begin{enumerate}[a)]
%\item
a)
$\max \{x,y\} = \dfrac{x+y+\sabs{x-y}}{2}$
%\item
\qquad
b)
$\min \{x,y\} = \dfrac{x+y-\sabs{x-y}}{2}$
%\end{enumerate}
\end{exercise}

\begin{exercise}
Encontre um número $M$ tal que $\sabs{x^3-x^2+8x} \leq M$ para todo $-2 \leq x \leq
10$.
\end{exercise}

\begin{exercise} \label{exercise:finishfuncsupinf}
Complete a demonstração de \propref{prop:funcsupinf}.
Isto é,
dado um conjunto $D$,
e duas funções limitadas
$f \colon D \to \R$ e $g \colon D \to \R$
tais que $f(x) \leq g(x)$ para todo $x \in D$, demonstre que
\begin{equation*}
\inf_{x\in D} f(x) \leq \inf_{x\in D} g(x) .
\end{equation*}
\end{exercise}

\begin{exercise} \label{exercise:supinffuncineq}
Sejam
$f \colon D \to \R$ e $g \colon D \to \R$ funções ($D$ não vazio).
\begin{enumerate}[a)]
%
\item
Suponha que
$f(x) \leq g(y)$ para todo $x \in D$ e todo $y \in D$. Mostre que
\begin{equation*}
\sup_{x\in D} f(x) \leq \inf_{x\in D} g(x) .
\end{equation*}
%
\item
Encontre conjuntos específicos $D$, $f$ e $g$ tais que
$f(x) \leq g(x)$ para todo $x \in D$, mas
\begin{equation*}
\sup_{x\in D} f(x) > \inf_{x\in D} g(x) .
\end{equation*}
\end{enumerate}
\end{exercise}

\begin{exercise}
Demonstre \propref{prop:funcsupinf} sem a hipótese de que
as funções são limitadas. Dica: é preciso usar os números reais
estendidos.
\end{exercise}

\begin{exercise} \label{exercise:sumofsup}
\pagebreak[2]
Seja $D$ um conjunto não vazio.
Suponha que $f \colon D \to \R$ e $g \colon D \to \R$ sejam funções limitadas.
\begin{enumerate}[a)]
\item
Mostre que
\begin{equation*}
%mbxlatex \begin{aligned}
%mbxlatex &
\sup_{x\in D} \bigl(f(x) + g(x) \bigr) \leq
\sup_{x\in D} f(x)
+
\sup_{x\in D} g(x)
\qquad \text{e}
%mbxSTARTIGNORE
\qquad
%mbxENDIGNORE
%mbxlatex \\
%mbxlatex &
\inf_{x\in D} \bigl(f(x) + g(x) \bigr) \geq
\inf_{x\in D} f(x)
+
\inf_{x\in D} g(x) .
%mbxlatex \end{aligned}
\end{equation*}
\item
Encontre um ou mais exemplos em que ocorram desigualdades estritas.
\end{enumerate}
\end{exercise}

\begin{exercise}
\pagebreak[2]
Suponha que $D$ seja não vazio, que $f \colon D \to \R$ e $g \colon D \to \R$ sejam
funções limitadas e que
$\alpha \in \R$.
\begin{enumerate}[a)]
\item
Mostre que $\alpha f \colon D \to \R$, definida por $(\alpha f) (x) \coloneqq \alpha
f(x)$, é uma função limitada.
\item
Mostre que $f+g \colon D \to \R$, definida por $(f+g) (x) \coloneqq f(x) + g(x)$,
é uma função limitada.
\end{enumerate}
\end{exercise}

\begin{exercise}
Sejam $f \colon D \to \R$ e $g \colon D \to \R$ funções, com $D$
não vazio e $\alpha \in \R$. Lembre-se do significado de $f+g$ e $\alpha f$ dado
no exercício anterior.
\begin{enumerate}[a)]
\item
Demonstre que, se $f+g$ e $g$ são limitadas, então $f$ é limitada.
\item
Encontre um exemplo em que $f$ e $g$ sejam ambas ilimitadas, mas $f+g$ seja
limitada.
\item
Demonstre que, se $f$ é limitada, mas $g$ é ilimitada, então $f+g$ é
ilimitada.
\item
Encontre um exemplo em que $f$ seja ilimitada, mas $\alpha f$ seja limitada.
\end{enumerate}
\end{exercise}

%%%%%%%%%%%%%%%%%%%%%%%%%%%%%%%%%%%%%%%%%%%%%%%%%%%%%%%%%%%%%%%%%%%%%%%%%%%%%%

\sectionnewpage

\section{Intervalos e o tamanho de \texorpdfstring{$\R$}{R}}
\label{sec:intandsizeR}

\sectionnotes{0,5--1 aula (a demonstração de que $\R$ não é enumerável pode ser opcional)}

Você certamente já viu a notação de intervalos\index{intervalo},
mas vamos dar aqui uma
definição formal. Para $a, b \in \R$ tais que $a < b$, definimos
\glsadd{not:openbdinterval}%
\glsadd{not:closedbdinterval}%
\glsadd{not:halfopenbdinterval}%
\glsadd{not:openubdinterval}%
\glsadd{not:closedubdinterval}%
\begin{align*}
& [a,b] \coloneqq \{ x \in \R : a \leq x \leq b \}, \\
& (a,b) \coloneqq \{ x \in \R : a < x < b \}, \\
& (a,b] \coloneqq \{ x \in \R : a < x \leq b \}, \\
& [a,b) \coloneqq \{ x \in \R : a \leq x < b \} .
\end{align*}
O intervalo $[a,b]$ é chamado de \emph{\myindex{intervalo fechado}}
e $(a,b)$ é chamado de \emph{\myindex{intervalo aberto}}. Os intervalos
da forma $(a,b]$ e $[a,b)$ são chamados de
\emph{intervalos semiabertos}\index{intervalo semiaberto}.

Os intervalos acima são \emph{intervalos limitados}\index{intervalo
limitado}, pois $a$ e $b$ são números reais.\footnote{Quando escrevemos
$(a,b)$, $[a,b]$ etc., sempre queremos dizer um intervalo limitado, a menos que digamos explicitamente
que $a$ ou $b$ pode ser $-\infty$ ou $\infty$.}
Definimos os \emph{intervalos não limitados}\index{intervalo não limitado},
\begin{align*}
& [a,\infty) \coloneqq \{ x \in \R : a \leq x \}, \\
& (a,\infty) \coloneqq \{ x \in \R : a < x \}, \\
& (-\infty,b] \coloneqq \{ x \in \R : x \leq b \}, \\
& (-\infty,b) \coloneqq \{ x \in \R : x < b \} .
\end{align*}
Para completar, definimos $(-\infty,\infty) \coloneqq \R$.
Às vezes, os intervalos $[a,\infty)$, $(-\infty,b]$ e $\R$ são chamados de
\emph{\myindex{intervalos fechados não limitados}},
e $(a,\infty)$, $(-\infty,b)$ e $\R$ são chamados de
\emph{\myindex{intervalos abertos não limitados}}.

Deixamos a demonstração da proposição a seguir como exercício.
Em resumo, um intervalo é um conjunto com pelo menos dois pontos
que contém todos os pontos entre quaisquer dois
de seus pontos.\footnote{Às vezes, conjuntos unitários e o conjunto vazio
também são chamados de intervalos, mas neste livro os intervalos têm pelo menos dois pontos.
Isto é, só definimos intervalos limitados quando $a < b$.}

\begin{prop} \label{prop:intervaldef}
Um conjunto $I \subset \R$ é um intervalo se, e somente se,
$I$ contém pelo menos dois pontos e
para todos $a,c \in I$ e $b \in \R$ tais que $a < b < c$, temos $b \in I$.
\end{prop}

Já vimos que todo intervalo aberto $(a,b)$ (com $a < b$, é claro)
é não vazio. Por exemplo, ele contém o número $\frac{a+b}{2}$.
Um fato surpreendente é que, do ponto de vista da teoria dos conjuntos,
todos os intervalos têm o mesmo \myquote{tamanho,} isto é, todos têm
a mesma cardinalidade. Por exemplo, a função $f(x) \coloneqq 2x$
leva o intervalo $[0,1]$ bijetivamente ao intervalo $[0,2]$.

Talvez ainda mais interessante seja a função $f(x) \coloneqq \tan(x)$,
que é uma função bijetiva de $(-\nicefrac{\pi}{2},\nicefrac{\pi}{2})$ em $\R$. Portanto, o intervalo
limitado $(-\nicefrac{\pi}{2},\nicefrac{\pi}{2})$ tem a mesma cardinalidade que $\R$. Não é
simples construir uma função bijetiva de $[0,1]$ para
$(0,1)$, mas isso é possível.

E não se preocupe: existe uma maneira de medir o \myquote{tamanho} dos subconjuntos
dos números reais que
\myquote{percebe}
a diferença entre $[0,1]$ e $[0,2]$. No entanto, sua definição adequada
requer muito mais ferramentas do que temos neste momento.

Digamos mais sobre a cardinalidade dos intervalos e, portanto, sobre a
cardinalidade de $\R$. Já vimos que existem números irracionais;
isto é,
$\R \setminus \Q$ não é vazio. A pergunta é: quantos números irracionais
existem? Acontece que há muito mais números irracionais do que racionais.
Vimos que $\Q$ é enumerável, e mostraremos
que $\R$ não é enumerável.
Na verdade, a cardinalidade de $\R$ é
igual à cardinalidade de $\sP(\N)$, embora não demonstremos
essa afirmação aqui.

\begin{thm}[Cantor]\index{teorema de Cantor}
$\R$ não é enumerável.
\end{thm}

Damos uma versão da demonstração original de Cantor,
de 1874, pois ela requer o mínimo de preparação. Normalmente, essa demonstração é apresentada
por contradição. No entanto, é mais fácil
entendê-la como uma demonstração construtiva direta de que qualquer subconjunto infinito enumerável
deve deixar de fora algum elemento.

\begin{proof}
Seja $X \subset \R$ um subconjunto infinito enumerável
tal que, para todo par de números reais
$a < b$, existe $x \in X$ tal que $a < x < b$. Se $\R$ fosse enumerável,
poderíamos tomar $X = \R$. Mostraremos que $X$ é necessariamente
um subconjunto próprio e, portanto, $X$ não pode ser igual a $\R$; logo, $\R$ não é
enumerável.

Como $X$ é infinito enumerável,
existe uma bijeção de $\N$ em $X$. Escrevemos $X$ como
uma sequência de números reais $x_1, x_2, x_3,\ldots$, tal que
cada número de $X$
é $x_n$ para algum $n \in \N$.

Construiremos indutivamente
duas sequências de números reais $a_1,a_2,a_3,\ldots$ e
$b_1,b_2,b_3,\ldots$. Defina
$a_1 \coloneqq x_1$ e $b_1 \coloneqq x_1+1$.
Observe que $a_1 < b_1$ e
$x_1 \notin (a_1,b_1)$.
Para algum $k > 1$,
suponha que $a_1,a_2,\ldots,a_{k-1}$ e $b_1,b_2,\ldots,b_{k-1}$ já tenham sido
definidos,
que valha
$a_1 < a_2 < \cdots < a_{k-1} < b_{k-1} < \cdots < b_2 < b_1$,
e que, para
cada $j=1,2,\ldots,k-1$,
tenhamos $x_\ell \notin (a_{j},b_{j})$
para $\ell=1,2,\ldots,j$.
\begin{enumerate}[(i)]
\item Defina $a_k \coloneqq x_n$, onde $n$ é o menor $n \in \N$
tal que $x_n \in (a_{k-1},b_{k-1})$. Tal $x_n$ existe pela nossa
hipótese sobre $X$, e $n \geq k$ pela hipótese sobre $(a_{k-1},b_{k-1})$.
\item Em seguida, defina $b_k$ como um número real em $(a_{k},b_{k-1})$,
por exemplo, $b_k \coloneqq \frac{a_k+b_{k-1}}{2}$.
\end{enumerate}
Observe que %$a_k < b_k$ e
$a_{k-1} < a_k < b_k < b_{k-1}$.
Observe também que $(a_{k},b_{k})$ não contém $x_k$ e, portanto,
não contém $x_{\ell}$ para $\ell=1,2,\ldots,k$.
Agora as duas sequências estão definidas.

Afirmação: \emph{$a_n < b_m$ para todos $n$ e $m$ em $\N$.} Demonstração: suponha primeiro
$n < m$. Então $a_n < a_{n+1} < \cdots < a_{m-1} < a_m < b_m$.
O mesmo vale para $n > m$. A afirmação está demonstrada.

Sejam $A \coloneqq \{ a_n : n \in \N \}$ e $B \coloneqq \{ b_n : n \in \N \}$.
Pela \propref{infsupineq:prop} e pela afirmação acima,
\begin{equation*}
\sup\, A \leq \inf\, B .
\end{equation*}
Defina $y \coloneqq \sup\, A$. O número $y$ não pode pertencer a $A$: se $y = a_n$
para algum $n$, então $y < a_{n+1}$, o que é impossível.
Da mesma forma, $y$ não pode pertencer a $B$. Portanto,
$a_n < y$ para todo $n\in \N$
e $y < b_n$ para todo $n\in \N$.
Em outras palavras, para todo $n \in \N$, temos $y \in (a_n,b_n)$.
Pela construção da sequência, $x_n \notin (a_n,b_n)$,
e, portanto, $y \neq x_n$. Como $y \neq x_n$ para todo $n \in \N$, temos $y
\notin X$.

Construímos um número real $y$ que não pertence a $X$, portanto $X$
é um subconjunto próprio de $\R$.
A sequência $x_1,x_2,\ldots$ não pode conter todos os elementos de $\R$,
e, portanto, $\R$ não é enumerável.
\end{proof}

\subsection{Exercícios}

\begin{exercise}
Para $a < b$, construa uma função bijetiva explícita de $(a,b]$ para $(0,1]$.
\end{exercise}

\begin{exercise}
Suponha que $f \colon [0,1] \to (0,1)$ seja uma função bijetiva.
Usando $f$, construa uma
função bijetiva de $[-1,1]$ para $\R$.
\end{exercise}

\begin{exercise} \label{exercise:intervaldef}
Demonstre \propref{prop:intervaldef}.
Isto é,
suponha que $I \subset \R$ seja um subconjunto com pelo menos dois elementos
tal que, se $a < b < c$ e $a, c \in I$, então $b \in I$.
Demonstre que $I$ é um dos nove tipos de intervalos explicitamente
definidos nesta seção.
Além disso, demonstre que todos os intervalos definidos nesta seção
satisfazem essa propriedade.
\end{exercise}

\begin{exercise}[Difícil]
Construa uma função bijetiva explícita de $(0,1]$ para $(0,1)$.
Dica: uma abordagem possível é a seguinte: primeiro leve $(\nicefrac{1}{2},1]$
a $(0,\nicefrac{1}{2}]$, depois leve
$(\nicefrac{1}{4},\nicefrac{1}{2}]$
a $(\nicefrac{1}{2},\nicefrac{3}{4}]$, e assim por diante.
Escreva a função explicitamente, isto é,
escreva um algoritmo que indique exatamente para onde cada número vai.
Em seguida, demonstre que a função é bijetiva.
\end{exercise}

\begin{exercise}[Difícil]
Construa uma função bijetiva explícita de $[0,1]$ para $(0,1)$.
\end{exercise}

\begin{exercise}
\leavevmode
\begin{enumerate}[a)]
\item
Mostre que todo intervalo fechado $[a,b]$ é a interseção
de uma quantidade enumerável de intervalos abertos.
\item
Mostre que todo intervalo aberto $(a,b)$
é uma união enumerável de intervalos fechados.
\item
Mostre que a interseção
de uma família não vazia e possivelmente infinita de intervalos fechados limitados,
$\bigcap\limits_{\lambda \in I} [a_\lambda,b_\lambda]$,
é vazia, um único ponto
ou um intervalo fechado limitado.
\end{enumerate}
\end{exercise}

\begin{exercise}
Suponha que $S$ seja um conjunto de intervalos abertos disjuntos em $\R$. Isto é,
se $(a,b) \in S$ e $(c,d) \in S$, então ou $(a,b) = (c,d)$
ou $(a,b) \cap (c,d) = \emptyset$. Demonstre que $S$ é enumerável.
\end{exercise}

\begin{exercise}
Demonstre que a cardinalidade de $[0,1]$ é igual à cardinalidade de
$(0,1)$ mostrando que
$\babs{[0,1]} \leq \babs{(0,1)}$ e
$\babs{(0,1)} \leq \babs{[0,1]}$.
Veja \defnref{def:comparecards}.
Esta demonstração requer o teorema de Cantor--Bernstein--Schr\"oder, que
enunciamos sem demonstração. Observe que esta demonstração não fornece
uma função bijetiva explícita.
\end{exercise}

\begin{exercise}[Desafiador]
Um número $x$ é \emph{algébrico}\index{número algébrico} se $x$ for raiz de
um polinômio não nulo com
coeficientes inteiros, ou seja, se $a_n x^n + a_{n-1} x^{n-1} + \cdots
+ a_1 x + a_0 = 0$, onde $a_0,a_1,\ldots,a_n \in \Z$, mas nem todos são zero.
\begin{enumerate}[a)]
\item
Mostre que há apenas
uma quantidade enumerável de números algébricos.
\item
Mostre que existem números não algébricos (transcendentes)
(siga os passos de Cantor e use a não enumerabilidade de $\R$).
\end{enumerate}
Dica: você pode usar o fato de que um polinômio não nulo de grau $n$ tem no máximo $n$ raízes
reais.
\end{exercise}

\begin{exercise}[Desafiador]
Seja $F$ o conjunto de todas as funções $f \colon \R \to \R$.
Demonstre que $\sabs{\R} < \sabs{F}$
usando o \thmref{cantorspowersetthm} de Cantor.\footnote{Curiosamente,
se $C$ é o conjunto das funções contínuas, então $\sabs{\R} = \sabs{C}$.}
\end{exercise}

%%%%%%%%%%%%%%%%%%%%%%%%%%%%%%%%%%%%%%%%%%%%%%%%%%%%%%%%%%%%%%%%%%%%%%%%%%%%%%

\sectionnewpage

\section{Representação decimal dos reais}
\label{sec:decimals}

\sectionnotes{1 aula (opcional)}

Costumamos pensar nos números reais por meio de suas
\emph{representações decimais}\index{representação decimal}.
Chamamos de (decimal)
\emph{\myindex{algarismo}}\index{algarismo decimal} um inteiro
entre $0$ e $9$.
Dado um inteiro positivo $n$, encontramos algarismos $d_K,d_{K-1},\ldots,d_2,d_1,d_0$ para algum
$K$ (cada $d_j$ é um inteiro entre $0$ e $9$) tais que
\begin{equation*}
n = d_K {10}^K + d_{K-1} {10}^{K-1} + \cdots + d_2 {10}^2 + d_1 10 + d_0 .
\end{equation*}
Frequentemente supomos que $d_K \neq 0$ (para evitar zeros à esquerda).
Para representar $n$, escrevemos a sequência de
algarismos: $n = d_K d_{K-1} \cdots d_2 d_1 d_0$.

De modo semelhante, representamos alguns números racionais. Para certos
números $x$, podemos encontrar
inteiros não negativos $K$ e $M$ e algarismos
$d_K,d_{K-1},\ldots,d_1,d_0,d_{-1},\ldots,d_{-M}$ tais que
\begin{equation*}
%mbxlatex \begin{aligned}
x
%mbxlatex &
= d_K {10}^K + d_{K-1} {10}^{K-1} + \cdots + d_2 {10}^2 + d_1 10 + d_0
%mbxlatex \\
%mbxlatex & \phantom{={}}
+ d_{-1} {10}^{-1} + d_{-2} {10}^{-2} + \cdots + d_{-M} {10}^{-M} .
%mbxlatex \end{aligned}
\end{equation*}
Escrevemos $x = d_K d_{K-1} \cdots d_1 d_0 \, . \, d_{-1} d_{-2} \cdots d_{-M}$.

Nem todo número real admite tal representação.
Nem mesmo o simples
número racional $\nicefrac{1}{3}$ a admite. O número irracional $\sqrt{2}$
também não admite tal representação. Para obter uma representação para
todos os números reais, precisamos permitir infinitos algarismos.

Consideremos apenas os números reais do intervalo $[0,1]$. Se
encontrarmos representações para eles, somando
números inteiros obtemos representações para todos os números reais.
Tomemos uma sequência infinita de algarismos decimais:
\begin{equation*}
0.d_1d_2d_3\ldots.
\end{equation*}
Temos um algarismo $d_j$ para cada $j \in \N$.
Renumeramos os algarismos para evitar os sinais negativos.
Escrevemos
\begin{equation*}
D_n \coloneqq
\frac{d_1}{10} +
\frac{d_2}{{10}^2} +
\frac{d_3}{{10}^3} +
\cdots +
\frac{d_n}{{10}^n} .
\end{equation*}
Dizemos que a
sequência de algarismos representa um número real $x$ se
\begin{equation*}
x =
\sup_{n \in \N} \left(
\frac{d_1}{10} +
\frac{d_2}{{10}^2} +
\frac{d_3}{{10}^3} +
\cdots +
\frac{d_n}{{10}^n}
\right) =
\sup_{n \in \N} \, D_n .
\end{equation*}
O número $D_n=0.d_1 d_2 d_3 \ldots d_n$ é o truncamento de $x$
a $n$ algarismos decimais.

\begin{prop} \label{prop:decimalprop}
\leavevmode
\begin{enumerate}[(i)]
\item
Toda sequência infinita de algarismos
$0.d_1d_2d_3\ldots$ representa um único número real $x \in [0,1]$, e
\begin{equation*}
D_n \leq x \leq D_n+\frac{1}{{10}^n} \qquad \text{para todo } n \in \N.
\end{equation*}
\item
Para todo $x \in (0,1]$ existe uma sequência infinita de algarismos
$0.d_1d_2d_3\ldots$ que representa $x$.
Existe uma única representação tal que
\begin{equation*}
D_n < x \leq D_n+\frac{1}{{10}^n} \qquad \text{para todo } n \in \N.
\end{equation*}
\end{enumerate}
\end{prop}

\begin{proof}
Comecemos pelo primeiro item. Tomemos uma sequência infinita arbitrária de
algarismos $0.d_1d_2d_3\ldots$. Usando a fórmula da soma geométrica, escrevemos
\begin{equation*}
\begin{split}
D_n =
\frac{d_1}{10} +
\frac{d_2}{{10}^2} +
\frac{d_3}{{10}^3} +
\cdots +
\frac{d_n}{{10}^n}
& \leq
\frac{9}{10} +
\frac{9}{{10}^2} +
\frac{9}{{10}^3} +
\cdots +
\frac{9}{{10}^n}
\\
& =
\frac{9}{10}
\bigl( 1 + \nicefrac{1}{10} +
{(\nicefrac{1}{10})}^2 + \cdots +
{(\nicefrac{1}{10})}^{n-1} \bigr)
\\
& =
\frac{9}{10}
\left(
\frac{1-{(\nicefrac{1}{10})}^{n}}{1-\nicefrac{1}{10}}
\right)
= 1-{(\nicefrac{1}{10})}^{n}
< 1 .
\end{split}
\end{equation*}
Em particular, $D_n < 1$ para todo $n$. Uma soma de números não negativos é
não negativa, logo $D_n \geq 0$, e portanto
\begin{equation*}
0 \leq \sup_{n\in \N} \, D_n \leq 1 .
\end{equation*}
Assim, $0.d_1d_2d_3\ldots$ representa um único número
$x \coloneqq \sup_{n\in \N} D_n \in [0,1]$.
Como $x$ é um supremo, $D_n \leq x$.
Tomemos $m \in \N$. Se $m < n$, então $D_m - D_n \leq 0$. Se $m > n$, então,
calculando como acima,
\begin{equation*}
D_m - D_n
=
\frac{d_{n+1}}{{10}^{n+1}} +
\frac{d_{n+2}}{{10}^{n+2}} +
\frac{d_{n+3}}{{10}^{n+3}} +
\cdots +
\frac{d_{m}}{{10}^m}
\leq
\frac{1}{{10}^{n}}
\bigl(
1-{(\nicefrac{1}{10})}^{m-n}
\bigr)
<
\frac{1}{{10}^{n}} .
\end{equation*}
Tomando o supremo em relação a $m$, obtemos
\begin{equation*}
x - D_n
\leq
\frac{1}{{10}^{n}} .
\end{equation*}

Passemos ao
segundo item. Tomemos qualquer $x \in (0,1]$.
Primeiro, provaremos a existência.
Por conveniência, seja $D_0 \coloneqq 0$.
Então,
$D_0 < x \leq D_0 + {10}^{-0}$.
Suponhamos que tenhamos definido os algarismos $d_1,d_2,\ldots,d_n$
e que
$D_k < x \leq D_k + {10}^{-k}$, para $k=0,1,2,\ldots,n$. Precisamos definir $d_{n+1}$.

Pela
\hyperref[thm:arch:i]{propriedade arquimediana} dos números reais,
encontramos um inteiro $j$ tal que
$x-D_n \leq j {10}^{-(n+1)}$. Tomemos o menor desses $j$ e obtemos
\begin{equation} \label{eq:theDnjineq}
(j-1){10}^{-(n+1)} < x-D_n \leq j {10}^{-(n+1)} .
\end{equation}
Seja $d_{n+1} \coloneqq j-1$.
Como $D_n < x$,
temos $d_{n+1} = j-1 \geq 0$. Por outro lado, como
$x-D_n \leq {10}^{-n}$, temos que
$j$ é no máximo 10 e, portanto, $d_{n+1} \leq 9$.
Logo, $d_{n+1}$ é um algarismo decimal.
Como $D_{n+1} = D_n + d_{n+1} {10}^{-(n+1)}$,
somemos $D_n$ à desigualdade
\eqref{eq:theDnjineq} acima:
\begin{multline*}
D_{n+1} = D_n + (j-1){10}^{-(n+1)} < x
\leq
D_n + j {10}^{-(n+1)} \\
=
D_n + (j-1) {10}^{-(n+1)} +
{10}^{-(n+1)} = D_{n+1} + {10}^{-(n+1)} .
\end{multline*}
Portanto, vale
$D_{n+1} < x \leq D_{n+1} + {10}^{-(n+1)}$.
Definimos indutivamente
uma sequência infinita de algarismos $0.d_1d_2d_3\ldots$.

Consideremos $D_{n} < x \leq D_{n} + {10}^{-n}$.
Como $D_n < x$ para todo $n$, temos
$\sup \{ D_n : n \in \N \} \leq x$.
A segunda desigualdade para $D_n$ implica
\begin{equation*}
x - \sup \{ D_m : m \in \N \}
\leq
x - D_n \leq 10^{-n} .
\end{equation*}
Como a desigualdade vale para todo $n$ e
${10}^{-n}$ pode ser arbitrariamente pequeno
(veja \exerciseref{exercise:bnlimit}),
temos $x \leq \sup \{ D_m : m \in \N \}$.
Portanto, $\sup \{ D_m : m \in \N \} = x$.

Resta provar a unicidade.
Suponhamos que $0.e_1e_2e_3\ldots$ seja outra representação de $x$.
Seja $E_n$ o truncamento com $n$ algarismos de $0.e_1e_2e_3\ldots$, e suponhamos
que $E_n < x \leq E_n + {10}^{-n}$ para todo $n \in \N$.
Suponhamos que, para algum $K \in \N$, $e_n = d_n$ para todo $n < K$; assim,
$D_{K-1} = E_{K-1}$. Então,
\begin{equation*}
E_K = D_{K-1} + e_K{10}^{-K} < x \leq E_K + {10}^{-K} = D_{K-1} +
e_K{10}^{-K} + {10}^{-K} .
\end{equation*}
Subtraindo $D_{K-1}$ e multiplicando por ${10}^{K}$, obtemos
\begin{equation*}
e_K < (x - D_{K-1}){10}^K \leq e_K + 1 .
\end{equation*}
De modo semelhante,
\begin{equation*}
d_K < (x - D_{K-1}){10}^K \leq d_K + 1 .
\end{equation*}
Portanto, tanto $e_K$ quanto $d_K$ são o maior inteiro $j$
tal que $j < (x - D_{K-1}){10}^K$, e, consequentemente, $e_K = d_K$.
Em outras palavras, a representação é única.
\end{proof}

A representação não é única se não exigirmos
$D_n < x$ para todo $n$.
Por exemplo, para o
número $\nicefrac{1}{2}$, o método usado na demonstração fornece a representação
\begin{equation*}
0.49999\ldots .
\end{equation*}
No entanto, $\nicefrac{1}{2}$ também tem a representação $0.50000\ldots$.

Os únicos números que têm representações não únicas são aqueles que terminam
em uma sequência infinita de algarismos $0$
ou em uma sequência infinita de algarismos $9$, pois a única representação para a qual
$D_n = x$ é aquela em que todos os algarismos após o $n$-ésimo são zeros.
Nesse caso,
há exatamente duas representações de $x$ (veja os exercícios).

Daremos outra demonstração de que os reais não são enumeráveis, usando
representações decimais.
Esta é a segunda demonstração de Cantor, provavelmente a mais conhecida.
Ela pode parecer mais curta, mas isso ocorre porque já fizemos
a parte mais difícil e nos resta um truque engenhoso para provar que $\R$ é
não enumerável. Esse truque é chamado de
\emph{\myindex{diagonalização de Cantor}}\index{diagonalização} e
também é usado em outras demonstrações.

\begin{thm}[Cantor]
O conjunto $(0,1]$ não é enumerável.
\end{thm}

\begin{proof}
Seja $X \coloneqq \{ x_1,x_2,x_3,\ldots \}$ um subconjunto enumerável qualquer dos números reais em $(0,1]$.
Construiremos um número real que não pertence a $X$. Seja
\begin{equation*}
x_n = 0.d_1^nd_2^nd_3^n\ldots
\end{equation*}
a representação única da proposição, isto é, $d_j^n$ é o
$j$-ésimo algarismo do $n$-ésimo número. Definamos
\begin{equation*}
e_n \coloneqq
\begin{cases}
1 & \text{se } d_n^n \neq 1, \\
2 & \text{se } d_n^n = 1.
\end{cases}
\end{equation*}
Seja $E_n$ o truncamento com $n$ algarismos de $y = 0.e_1e_2e_3\ldots$. Como
todos os algarismos são não nulos, $E_n < E_{n+1} \leq y$. Portanto,
\begin{equation*}
E_n < y \leq E_n + {10}^{-n}
\end{equation*}
para todo $n$, e, pela proposição,
essa é a representação única de $y$. Para todo $n$, o $n$-ésimo algarismo
de $y$ é diferente do $n$-ésimo algarismo de $x_n$, logo $y \neq x_n$.
Assim, $y \notin X$. Como $X$ era um subconjunto enumerável arbitrário,
$(0,1]$ não é enumerável. Veja um
exemplo na \figureref{cantorexamplefig}.
\begin{myfigureht}
\myincludepdft{cantorexamplefig}{%
Cinco números, indicados por x sub 1 até x sub 5, aparecem alinhados um
abaixo do outro na notação decimal.
Os cinco primeiros algarismos decimais de cada número estão visíveis.
Em particular, x sub 1 é 0 vírgula 1 3 2 1 0, seguido de reticências que
indicam a continuação dos algarismos.
O primeiro algarismo, 1, está contornado por uma caixa.
x sub 2 é 0 vírgula 7 9 4 1 3, seguido de reticências.
O segundo algarismo, 9, está contornado por uma caixa.
x sub 3 é 0 vírgula 3 0 1 3 4, seguido de reticências.
O terceiro algarismo, 1, está contornado por uma caixa.
x sub 4 é 0 vírgula 8 9 2 5 6, seguido de reticências.
O quarto algarismo, 5, está contornado por uma caixa.
x sub 5 é 0 vírgula 1 6 0 2 4, seguido de reticências.
O quinto algarismo, 4, está contornado por uma caixa.
Reticências adicionais indicam que a sequência de números continua.
À direita, aparece o texto: Número que não está na lista:
y igual a 0 vírgula 2 1 2 1 1, seguido de reticências.}
\caption{Exemplo da diagonalização de Cantor, com os algarismos diagonais $d_n^n$ marcados.\label{cantorexamplefig}}
\end{myfigureht}
\end{proof}

Usando algarismos decimais, também podemos encontrar muitos números que não são racionais.
A proposição a seguir vale para todo
número racional, mas, por simplicidade, enunciaremos o resultado apenas para $x \in (0,1]$.

\begin{prop} \label{prop:rationaldecimal}
Se $x \in (0,1]$ é racional e $x = 0.d_1d_2d_3\ldots$,
então os algarismos decimais tornam-se periódicos a partir de algum ponto. Isto é, existem
inteiros positivos $N$ e $P$ tais que, para todo $n \geq N$, $d_n = d_{n+P}$.
\end{prop}

\begin{proof}
Suponhamos que $x = \nicefrac{p}{q}$, em que $p$ e $q$ são inteiros positivos.
Se $x$ tiver uma representação decimal não única, então ambas as representações
terão algarismos periódicos a partir de algum ponto:
ou serão todos $0$, ou todos $9$;
veja \exerciseref{exercise:nonuniquedecimals}.
Portanto, suponhamos que $x$ tenha uma representação única.
Também suporemos que $x \neq 1$, de modo que $p < q$.

Para calcular o primeiro algarismo, dividimos $10p$ por
$q$. Seja $d_1$ o quociente.
O resto $r_1$ é um inteiro
entre $0$ e $q-1$. Isto é, $d_1$ é o maior inteiro
tal que $d_1 q \leq 10p$ e, em seguida, $r_1 = 10p - d_1q$.
Como $p < q$, temos $d_1 < 10$, logo $d_1$ é um algarismo.
Além disso,
\begin{equation*}
\frac{d_1}{10} \leq \frac{p}{q} =
\frac{d_1}{10} + \frac{r_1}{10q} \leq \frac{d_1}{10} +
\frac{1}{10} .
\end{equation*}
A primeira desigualdade é estrita, pois $x$ tem uma
representação única. Portanto, $d_1$ é de fato o primeiro algarismo.
Resta $\frac{r_1}{10q}$.
Para calcular $d_2$, dividimos $10r_1$ por $q$, obtendo $d_2$
e $r_2$, em que $r_2 = 10r_1 - d_2 q$.
Isolando $r_1$ e substituindo em
$\frac{p}{q} = \frac{d_1}{10} + \frac{r_1}{10 q}$,
obtemos
$\frac{p}{q} = \frac{d_1}{10} + \frac{d_2}{{10}^2} +
\frac{r_2}{{10}^2 q} = D_2 +
\frac{r_2}{{10}^2 q}$.
Como $r_2 < q$, novamente temos $d_2 < 10$.
Para calcular $d_3$, dividimos $10r_2$ por $q$,
obtendo $d_3$ e $r_3$. E assim por diante.
Depois de calcular $n-1$ algarismos, temos
$\frac{p}{q} = D_{n-1} + \frac{r_{n-1}}{10^{n-1} q}$.
Para obter o $n$-ésimo algarismo,
dividimos $10r_{n-1}$ por $q$
para obter o quociente $d_n$, o resto $r_n$ e as desigualdades
\begin{equation*}
\frac{d_n}{10} \leq \frac{r_{n-1}}{q} =
\frac{d_n}{10} + \frac{r_n}{10q} \leq \frac{d_n}{10} +
\frac{1}{10} .
\end{equation*}
Dividindo por $10^{n-1}$ e somando $D_{n-1}$, obtemos
\begin{equation*}
D_n \leq D_{n-1} + \frac{r_{n-1}}{10^{n-1} q} = \frac{p}{q} \leq D_n +
\frac{1}{10^n} .
\end{equation*}
Pela unicidade, o $n$-ésimo algarismo $d_n$ é o algarismo obtido pela construção.

O novo algarismo depende apenas do resto do passo anterior.
Há no máximo $q$ restos possíveis.
Portanto, o processo deve começar a se repetir após no máximo $q$ passos,
e, assim, $P$ é no máximo $q$.
\end{proof}

A recíproca da proposição também vale e fica como exercício.

\begin{example}
O número
\begin{equation*}
x = 0.101001000100001000001\ldots
\end{equation*}
é irracional.
Mais precisamente, a sequência de algarismos é 1,
depois 0, depois 1, depois dois zeros, depois 1,
depois três zeros,
e assim por diante.
A irracionalidade de $x$ segue da
proposição: os algarismos nunca se tornam periódicos.
Para todo $P$,
se avançarmos suficientemente, encontraremos um 1 seguido de pelo menos
$P$ zeros, e há infinitos algarismos 1 desse tipo.
Isto é, existem infinitos $n$ (logo, valores de $n$ arbitrariamente grandes)
tais que $d_n=1$ e $d_{n+P}=0$; em outras palavras,
$d_n \neq d_{n+P}$.
\end{example}

\subsection{Exercícios}

\begin{exercise}[Fácil]
Qual é a representação decimal de $1$ garantida por
\propref{prop:decimalprop}? Mostre que ela de fato satisfaz
a condição.
\end{exercise}

\begin{exercise}
Demonstre a recíproca de \propref{prop:rationaldecimal}, isto é,
se os algarismos da representação decimal de $x$ forem periódicos a partir de algum ponto, então
$x$ será racional.
\end{exercise}

\begin{exercise} \label{exercise:nonuniquedecimals}
Mostre que os números reais $x \in (0,1)$ com representação decimal não única
são exatamente os números racionais que podem ser escritos
como $\frac{m}{10^n}$ para alguns inteiros $m$ e $n$. Neste caso, mostre que
existem exatamente duas representações de $x$.
\end{exercise}

\begin{exercise} \label{exercise:decimalpropbaseb}
Seja $b \geq 2$ um inteiro. Defina uma representação de um número real em
$[0,1]$ na base $b$, em vez da base 10, e demonstre
\propref{prop:decimalprop} para a base $b$.
\end{exercise}

\begin{exercise}
Usando o exercício anterior com $b=2$ (sistema binário),
mostre que a cardinalidade de $\R$ é igual à cardinalidade de $\sP(\N)$,
obtendo assim mais uma demonstração (embora relacionada) de que $\R$ não é enumerável.
Dica: Construa duas injeções, uma de $[0,1]$ em $\sP(\N)$
e outra de $\sP(\N)$ em $[0,1]$. Dica 2: Dado um
conjunto $A \subset \N$, a função que faz o $n$-ésimo algarismo binário de $x$ ser 1 se
$n\in A$ funciona quase, mas não é injetiva. Modifique-a.
\end{exercise}

\begin{exercise}[Desafiador]
Construa explicitamente uma injeção de $[0,1] \times [0,1]$ em $[0,1]$ (pense por que
isso é tão surpreendente\footnote{%
Com um pouco mais de trabalho (ou aplicando o
teorema de Cantor--Bernstein--Schr\"oder),
é possível demonstrar que existe uma bijeção.
Ao provar esse resultado, Cantor aparentemente escreveu:
\myquote{Vejo isso, mas não acredito.}}).
Em seguida, descreva o conjunto dos números de $[0,1]$ que não pertencem à imagem da sua
injeção (a menos, é claro, que você tenha conseguido construir uma bijeção).
Dica: Considere os algarismos pares e ímpares da expansão decimal.
\end{exercise}

\begin{exercise}
Demonstre que, se $x = \nicefrac{p}{q} \in (0,1]$ é racional, $q > 1$,
então o período $P$ dos algarismos periódicos da representação decimal
de $x$ é, na verdade, menor ou igual a $q-1$.
\end{exercise}

\begin{exercise} \label{exercise:bnlimit}
Demonstre que, se $b \in \N$ e $b \geq 2$, então, para todo $\epsilon > 0$,
existe um $n \in \N$ tal que
$b^{-n} < \epsilon$. Dica:
Uma possibilidade é primeiro demonstrar por indução que $b^n > n$ para todo $n \in \N$.
\end{exercise}

\begin{exercise}
Construa explicitamente uma injeção $f \colon \R \to \R \setminus \Q$ usando
\propref{prop:rationaldecimal}.
\end{exercise}
